\documentclass{article} % For LaTeX2e
\usepackage{iclr2027_conference,times}

\input{math_commands.tex}

\usepackage{amsmath}
\usepackage{amssymb}
\usepackage{algorithm}
\usepackage{algpseudocode}
\usepackage{graphicx}
\usepackage{url}
% hyperref must be loaded AFTER amsmath, otherwise \eqref expands to
% "equation~N" instead of "(N)".  This was the cause of the "Eq. equation 10"
% strings in the previous build.
\usepackage{hyperref}

\usepackage{subcaption}
\captionsetup[subfigure]{font=small, labelfont=bf, skip=2pt}
\usepackage{booktabs}
\usepackage{multirow}
\usepackage{array}
\usepackage{xcolor}
\definecolor{bestblue}{HTML}{14479E}
\definecolor{secondorange}{HTML}{BF5B04}
\newcommand{\best}[1]{\textcolor{bestblue}{\textbf{#1}}}
\newcommand{\second}[1]{\textcolor{secondorange}{#1}}
\newcommand{\rankkey}{\best{best} and \second{second-best}}
\newcommand{\tbd}{\textcolor{red}{--}}
% The two reliability variants of the particle score R_m (Eq. 3), as points of the
% (beta, gamma, lambda_D) plane: /E = (1, 1.3, 0) prices context uncertainty through
% the entropy channel; /D = (1, 1, 0.3) adds trajectory disagreement on top of a full
% unit of that entropy penalty.  Everything else is identical.
\newcommand{\lorekvE}{LORE-KV\,/E}
\newcommand{\lorekvD}{LORE-KV\,/D}
\title{Sampling the Future: Monte-Carlo Query Estimation for KV Cache Eviction}
\author{Antiquus S.~Hippocampus, Natalia Cerebro \& Amelie P. Amygdale \thanks{ Use footnote for providing further information
about author (webpage, alternative address)---\emph{not} for acknowledging
funding agencies.  Funding acknowledgements go at the end of the paper.} \\
Department of Computer Science\\
Cranberry-Lemon University\\
Pittsburgh, PA 15213, USA \\
\texttt{\{hippo,brain,jen\}@cs.cranberry-lemon.edu} \\
\And
Ji Q. Ren \& Yevgeny LeNet \\
Department of Computational Neuroscience \\
University of the Witwatersrand \\
Joburg, South Africa \\
\texttt{\{robot,net\}@wits.ac.za} \\
\AND
Coauthor \\
Affiliation \\
Address \\
\texttt{email}
}

\newcommand{\fix}{\marginpar{FIX}}
\newcommand{\new}{\marginpar{NEW}} \usepackage[textsize=tiny]{todonotes}
 \newcommand{\todom}[1]{\todo[color=magenta!10, inline]{\small MV: #1}}


\begin{document}


\maketitle

\begin{abstract}

\todom{
 how about we say something more catchy and very understandable for aanywone:  in this paper, we move towardss asking “which memory will matter while answering?” rather than “which memory appeared important while reading?”
}


Most KV-cache eviction methods ask, in effect, \emph{``which memory appeared important while reading the prompt?''} We instead ask, \emph{``which memory will matter while answering?''} The difficulty is that future decoding queries are unknown when the cache is compressed. We address this by treating the future-query distribution as something to \emph{sample} rather than guess. We introduce \textbf{LORE-KV} (\emph{\textbf{\underline{L}}ookahead \textbf{\underline{O}}utput-perturbation with \textbf{\underline{R}}eliability-weighted \textbf{\underline{E}}nsembles for \textbf{\underline{K}}ey--\textbf{\underline{V}}alue caches}), a training-free KV-cache eviction method. LORE-KV samples $M$ short continuations from the frozen model, uses their query states as a Monte-Carlo estimate of future decoding queries, and scores each KV pair by its projected leave-one-out attention-output perturbation under this ensemble. Each trajectory receives a reliability weight based on branch atypicality, context uncertainty, and, optionally, disagreement with the other trajectories in query space. The lookaheads are discarded before final generation, requiring no training, auxiliary model, or modification of the base LLM.  \textcolor{red}{Across LongBench, RULER, and Needle-in-a-Haystack, LORE-KV remains consistently effective across models, tasks, and cache budgets, with its largest gains under aggressive compression; moreover, the gains saturate with only a few lookahead trajectories and increase with prompt-to-future query mismatch, showing that the improvement comes primarily from better future-query estimation rather than simply spending more test-time compute.}

\end{abstract}


\section{Introduction}
\label{sec:introduction}

Autoregressive large language models (LLMs) rely on key-value (KV) caching to reuse previously computed hidden states during decoding. While KV caching greatly accelerates generation, its memory and bandwidth cost grows with sequence length and batch size, making long-context inference increasingly expensive \citep{behnam2025rocketkv,li2025kvtuner,lancucki2025hyperscaling}. KV-cache eviction reduces this cost by retaining only the KV pairs expected to be useful during generation. Existing methods estimate token importance using recency or attention sinks \citep{xiao2023streamingllm}, accumulated or prompt-end attention \citep{zhang2023h2o,liu2023scissorhands,li2024snapkv}, key geometry \citep{devoto2024simple,park2026keydiff}, or learned retention signals \citep{moschella2026learning,bui2026make}.

A central limitation of existing KV-cache eviction methods is that they often estimate token importance using prompt-side queries, even though the retained cache is later used by future decoding queries. Lookahead Q-Cache \citep{wang2025lookahead} partially addresses this by using a single surrogate response, while learned methods such as LookaheadKV estimate future utility with trained predictors \citep{ahn2026lookaheadkv,moschella2026learning,bui2026make}; however, a single surrogate captures only one possible future, and learned approaches require extra training and model-specific parameters. At the same time, attention alone is not a complete measure of deletion impact, because tokens with similar attention can contribute differently through their value vectors \citep{guo2024criticalkv}. Output-aware methods such as AnDPro and LaProx improve the scoring criterion \citep{geng2026accurate,mai2026laprox}, but still generally rely on prompt-side queries. Thus, prior work improves either the future-query estimate or the deletion score, but not both together over multiple plausible future continuations.

We propose \textbf{\underline{LORE-KV}}, short for \emph{\textbf{\underline{L}}ookahead \textbf{\underline{O}}utput-Perturbation with \textbf{\underline{R}}eliability-weighted \textbf{\underline{E}}nsembles}, a training-free KV-cache eviction method that \emph{samples the future} instead of relying on prompt-side queries or a single surrogate continuation. LORE-KV generates multiple short continuations from the frozen model, uses their query states to approximate plausible future decoding queries, and scores each KV pair using projected leave-one-out output perturbation. \textcolor{blue}{It combines response-side lookahead from Lookahead Q-Cache \citep{wang2025lookahead} with projected output-aware scoring from AnDPro \citep{geng2026accurate}, extending both with multiple trajectories and reliability-weighted aggregation} Unlike learned future-utility methods \citep{ahn2026lookaheadkv,moschella2026learning,bui2026make,yang2026indexmem}, LORE-KV requires no offline training or auxiliary predictor, adds only a small one-time lookahead cost, and empirically gains most when prompt-side and future queries differ, supporting better query estimation rather than extra test-time compute as the main source of improvement.


\todom{The strongest novelty seems to be the combination of future-query estimation with projected deletion scoring. Could we state somwhere  explicitly that response-side lookahead is "inherited" from Lookahead Q-Cache and projected output-aware scoring from AnDPro, then isolate empirically what multi-trajectory confidence weighting adds beyond them?}


\textbf{Contributions.}
Our contributions are threefold. \textbf{First}, we identify and empirically characterize the \emph{query-reference mismatch} in KV-cache eviction, showing that prompt-side queries can be poor proxies for the future decoding queries that consume the retained cache. \textcolor{blue}{\textbf{Second}, we introduce LORE-KV, a training-free combination of multiple sampled future-query trajectories, projected leave-one-out deletion scoring, and reliability-weighted aggregation. Unlike learned future-utility methods, LORE-KV requires no offline training, auxiliary predictor, or model-specific learned parameters, while directly adapting its query estimate to each input at test time.} \textbf{Third}, we evaluate LORE-KV across LongBench, RULER, and Needle-in-a-Haystack on two LLM backbones and multiple cache budgets, with ablations showing that its gains primarily arise from improved future-query estimation rather than additional test-time compute.

\todom{ suggestion: somehow, what we do is: 
training-free combination:
multiple sampled future-query trajectories + projected leave-one-out scoring + confidence aggregation.

 at the same time, there were combinations that were trained, so naturally people will think, "Well, with training I can do better." I think we should only then defend it. What is the benefit of non-training: computational efficiency, collapse avoidance, or whatever, and back to it empirical evidence
 } 

\section{Methodology} \label{sec:methodology}
LORE-KV has three stages: (i) $M$ short continuations form a Monte-Carlo estimate of future decoding queries (Section~\ref{subsec:lookahead}); (ii) each trajectory receives a reliability weight based on likelihood, entropy, or ensemble disagreement (Section~\ref{subsec:confidence}); and (iii) a projected counterfactual deletion score measures the effect of evicting each KV pair under the weighted future queries (Sections~\ref{subsec:deletion}--\ref{subsec:core_score}). Algorithm~\ref{alg:lorekv} summarizes the method, Table~\ref{tab:hyperparams} lists the hyperparameters, and Appendices~\ref{app:method_details} and~\ref{sec:ablation_components} provide implementation details and ablations.

\subsection{Problem Setup} \label{subsec:problem_setup}
Let $f_\theta$ denote a frozen pretrained transformer language model, whose parameters $\theta$ remain fixed throughout. Given a prompt $X=(x_1,\ldots,x_n)$, a single prefill pass produces the cached keys and values $\mathbf{K}^{l,h}=\{\mathbf{k}^{l,h}_{1},\ldots,\mathbf{k}^{l,h}_{n}\}$ and $\mathbf{V}^{l,h}=\{\mathbf{v}^{l,h}_{1},\ldots,\mathbf{v}^{l,h}_{n}\}$ for each layer $l\in\{1,\ldots,L\}$ and KV head $h\in\{1,\ldots,H\}$, together with the prompt-window queries $\mathcal{Q}^{l,h}_{\mathcal{W}}$. Let $\mathcal{W}=\{n-w+1,\ldots,n\}$ denote the prompt-end observation window of size $w$, and let $\mathcal{B}$ denote the optional structural-token positions (more details in Appendix~\ref{app:setup_details}). We also retain the first prompt token, which commonly acts as an attention sink \citep{xiao2023streamingllm}. Thus, the protected and candidate sets are $\mathcal{P}=\{1\}\cup\mathcal{W}\cup\mathcal{B}$ and $\mathcal{C}=\{1,\ldots,n\}\setminus\mathcal{P}$, respectively, with only tokens in $\mathcal{C}$ scored for eviction.


\subsection{Test-Time Lookahead Query Construction}
\label{subsec:lookahead}

Prompt-side queries may poorly represent the response-side queries used during generation. To reduce this mismatch, LORE-KV uses a small, fixed test-time budget before final decoding to sample $M$ short lookahead continuations from the frozen model:
\begin{equation}
Y^{(m)}\sim f_\theta(\cdot\mid X),\qquad
Y^{(m)}=\{y^{(m)}_1,\ldots,y^{(m)}_T\},\qquad
m=1,\ldots,M,
\label{eq:lookahead_cont}
\end{equation}
where $T$ is small. The continuations are used as temporary probes for estimating future decoding queries. To keep the reliability score correctly centered, each trajectory is evaluated under the same distribution from which it is sampled (Section~\ref{subsec:confidence}). For each trajectory $m$, we store its lookahead query states as $ \mathcal{Q}^{(m),l,h}_{\mathrm{look}} = \left\{ \mathbf{q}^{(m),l,h}_{1}, \ldots, \mathbf{q}^{(m),l,h}_{T} \right\}, $ and combine them with the prompt-window queries for scoring:

\begin{equation}
\mathcal{Q}^{l,h}_{\mathrm{score}}
=
\mathcal{Q}^{l,h}_{\mathcal{W}}
\cup
\left(
\bigcup_{m=1}^{M}
\mathcal{Q}^{(m),l,h}_{\mathrm{look}}
\right).
\label{eq:score_queries}
\end{equation}

The $M$ lookahead trajectories therefore provide response-side query samples for subsequent KV scoring. Figure~\ref{fig:future_query_alignment} shows that $\mathcal{Q}^{l,h}_{\mathrm{score}}$ better covers the true future queries than $\mathcal{Q}^{l,h}_{\mathcal{W}}$ alone, while Appendix~\ref{app:lookahead_details} relates this construction to prior surrogate-query methods.

\subsection{Particle Reliability Weighting: Entropy and Trajectory Disagreement} \label{subsec:confidence}

The lookahead trajectories can differ in reliability, so LORE-KV assigns each a weight. Let $u\in{1,\ldots,T}$ index the $T$ lookahead steps, and let $h^{(m)}u=(X,y^{(m)}{<u})$ denote the context of trajectory $m$ at step $u$. Let $q(\cdot\mid h)$ be the sampling distribution. Raw log-probability mixes two effects: whether the sampled token is atypical for its local context and how uncertain that context is. We separate these channels with the centered score
\begin{equation}
 c^{(m)}_u
 =\log q\!\left(y^{(m)}_u\mid h^{(m)}_u\right)
 +H\!\left(q(\cdot\mid h^{(m)}_u)\right),
 \label{eq:centered_score}
\end{equation}
for which $\mathbb{E}[c^{(m)}_u\mid h^{(m)}_u]=0$. With $\bar c_m=\tfrac1T\sum_u c^{(m)}_u$ and $\bar H_m=\tfrac1T\sum_u H(q(\cdot\mid h^{(m)}_u))$, trajectory reliability is
\begin{equation}
 R_m=\beta\,\bar c_m-\gamma\,\bar H_m-\lambda_D\,\hat D_m,
 \label{eq:confidence}
\end{equation}
where $\bar c_m$ measures branch atypicality relative to the uncertainty of the visited contexts, $\bar H_m$ penalizes high-uncertainty histories, and $\hat D_m$ is standardized disagreement between the trajectory's lookahead query states and those of the other particles. The definition, standardization, calibration, and finite-$M$ edge cases of $\hat D_m$ are given in Appendix~\ref{app:reliability_details}. The reliability scores are converted to soft particle weights,
\begin{equation}
 \omega_m=\frac{\exp(R_m)}{\sum_{r=1}^{M}\exp(R_r)},
 \qquad
 \boldsymbol{\omega}=\operatorname{softmax}(\mathbf{R})\in\Delta^{M-1}.
 \label{eq:particle_weight}
\end{equation}
These weights control each trajectory's contribution to the future-query term in Section~\ref{subsec:core_score}. Soft weighting is used rather than selecting a single trajectory because the measured reliability signals are only imperfect proxies for oracle trajectory utility (Appendix~\ref{app:confidence_details}).

We consider two settings of Eq.~\eqref{eq:confidence}. \emph{\lorekvE} uses $(\beta,\gamma,\lambda_D)=(1,1.3,0)$ and is the default. \emph{\lorekvD} uses $(1,1,0.3)$ and additionally down-weights trajectories that are atypical in the query space used by the eviction score. The variants reuse the same trajectories and differ only in reliability coefficients, with equal decoding cost. Appendix~\ref{app:reliability_details} provides details, and Table~\ref{tab:ablation_summary} reports single-factor controls.

\subsection{Closed-Form Counterfactual Deletion Effect} \label{subsec:deletion}
With the scoring queries in hand, we now specify \emph{what} is measured for each candidate token. For a scoring query $\mathbf{q}^{l,h}$ and candidate token $i\in\mathcal{C}$, the full-cache attention weight is
$a^{l,h}_{q,i}=\operatorname{Softmax}_i\!\left((\mathbf{q}^{l,h})^\top\mathbf{k}^{l,h}_{i}/\sqrt{d}\right)$,
where $d$ is the head dimension, and the corresponding full-cache attention output is
$\mathbf{o}^{l,h}_{q}=\sum_{j=1}^{n}a^{l,h}_{q,j}\mathbf{v}^{l,h}_{j}$.
We write
$(a^{l,h}_{q,i},\mathbf{o}^{l,h}_{q})
=\operatorname{Attn}(\mathbf{q}^{l,h},\mathbf{K}^{l,h},\mathbf{V}^{l,h})$
for this full-cache attention weight--output pair. If token $i$ is removed and the remaining attention weights are renormalized, the post-deletion output becomes
$\mathbf{o}^{l,h}_{q,-i}
=\sum_{j\neq i}\frac{a^{l,h}_{q,j}}{1-a^{l,h}_{q,i}+\epsilon}\mathbf{v}^{l,h}_{j}
=\frac{\mathbf{o}^{l,h}_{q}-a^{l,h}_{q,i}\mathbf{v}^{l,h}_{i}}
{1-a^{l,h}_{q,i}+\epsilon}$,
where $\epsilon$ prevents numerical instability. The corresponding leave-one-out deletion effect is therefore
\begin{equation}
\Delta\mathbf{o}^{l,h}_{q,i}
=\mathbf{o}^{l,h}_{q}-\mathbf{o}^{l,h}_{q,-i}
\approx
\frac{a^{l,h}_{q,i}}{1-a^{l,h}_{q,i}+\epsilon}
\left(\mathbf{v}^{l,h}_{i}-\mathbf{o}^{l,h}_{q}\right).
\label{eq:deletion_effect}
\end{equation}
Eq.~\eqref{eq:deletion_effect} shows why attention alone is insufficient: a token is important only when it receives non-negligible attention and its value meaningfully differs from the current attention output.

\textbf{Output-projection-aware causal score.}
Let $\mathbf{W}^{O}_{l,h}$ denote the output-projection block for head $h$ in layer $l$. LORE-KV scores the residual-stream impact of deleting token $i$ as
\begin{equation}
e^{l,h}_{q,i}
=
\left\|
\mathbf{W}^{O}_{l,h}\Delta\mathbf{o}^{l,h}_{q,i}
\right\|_2^2 .
\label{eq:projected_score}
\end{equation}
Unlike raw attention, this score measures the deletion effect after projection into the model stream. Further details are given in Appendix~\ref{app:lookahead_details}. Figure~\ref{fig:cache_distribution} evaluates the resulting cache fidelity, while Figure~\ref{fig:layer_head_analysis} shows how the perturbation is distributed across layers and heads.

\subsection{Core Score, Selection, and Compressed Cache}
\label{subsec:core_score}

For each token $i$, we collect perturbation scores over the prompt-window queries, $\mathcal{E}^{l,h}_{i,\mathcal{W}} =\{e^{l,h}_{q,i}:q\in\mathcal{Q}^{l,h}_{\mathcal{W}}\}$, and each lookahead trajectory, $\mathcal{E}^{l,h}_{i,m} =\{e^{l,h}_{q,i}:q\in\mathcal{Q}^{(m),l,h}_{\mathrm{look}}\}$. Using mean aggregation, $\rho(\mathcal{E})=\operatorname{Mean}(\mathcal{E})$, with further details in Appendix~\ref{sec:ablation_components}, the final token score is

\begin{equation}
S^{l,h}_{i}
=\rho(\mathcal{E}^{l,h}_{i,\mathcal{W}})
+\lambda_L\sum_{m=1}^{M}\omega_m\,\rho(\mathcal{E}^{l,h}_{i,m}),
\label{eq:final_score}
\end{equation}
where $\lambda_L$ governs how much the future-query information counts, and we stack the per-token scores of a head into the vector $\mathbf{S}^{l,h}=(S^{l,h}_i)_{i\in\mathcal{C}}$. At $\lambda_L=0$, LORE-KV collapses to a prompt-window output-perturbation method; as $\lambda_L$ grows, the policy leans more heavily on the future queries. The lookahead term is not a cosmetic re-weighting of the prompt-side score: Figure~\ref{fig:score_comparison} shows that roughly one in five of the tokens retained at $B=128$ differs between the two rules, and Table~\ref{tab:ablation_summary} prices the $\lambda_L{=}0$ setting directly.

Two operators turn these scores into a compressed cache, and both reduce to their simplest form in the core method:
\begin{equation}
\{B_{l,h}\}_{h}=\operatorname{Alloc}\!\left(B_l,\{E_{l,h}\}_{h};\tau\right),
\qquad
\mathcal{R}^{l,h}=\operatorname{Select}\!\left(\mathbf{S}^{l,h},B_{l,h}\right).
\label{eq:alloc_select}
\end{equation}
For the core method, $\operatorname{Alloc}$ with $\tau=0$ returns the uniform split $B_{l,h}=\lfloor B_l/H\rfloor$ of the per-layer budget $B_l$ (reserved after protected tokens), and $\operatorname{Select}$ keeps the top-$B_{l,h}$ candidates by score. The final retained set and the compressed cache are
\begin{equation}
\mathcal{R}^{l,h}_{\mathrm{final}}=\mathcal{P}\cup\mathcal{R}^{l,h},
\qquad
\mathbf{K}^{l,h}_{\star},\mathbf{V}^{l,h}_{\star}
=\operatorname{Gather}\left(\mathbf{K}^{l,h},\mathbf{V}^{l,h},\mathcal{R}^{l,h}_{\mathrm{final}}\right),
\label{eq:gather}
\end{equation}
at which point the lookahead continuations are thrown away and final decoding runs on $(\mathbf{K}_\star,\mathbf{V}_\star)$. The switches $\tau>0$ and $\lambda_k,\lambda_v>0$ turn $\operatorname{Alloc}$ and $\operatorname{Select}$ into the optional plug-in rules of Appendix~\ref{sec:ablation_components}, which are off throughout the core method; each recovers its default exactly at zero (Table~\ref{tab:hyperparams}).

\section{Experiments} \label{subsec:setup}
We follow the evaluation protocol of AnDPro \citep{geng2026accurate}, comparing all methods at matched cache budgets, backbones, and observation windows. We evaluate on \textbf{LongBench} \citep{bai2024longbench}, \textbf{RULER} \citep{hsieh2024ruler}, and \textbf{Needle-in-a-Haystack} \citep{kamradt2023needle} using Mistral-7B-Instruct-v0.2 \citep{jiang2023mistral} and Llama-3.1-8B-Instruct \citep{dubey2024llama}. We compare against Full Cache and all KV-eviction baselines reported in Table~\ref{tab:longbench_mistral}. We sweep $B\in\{128,256,512,1024\}$ with observation window $w=32$ and chunk size $4$; Table~\ref{tab:longbench_datasets} lists the LongBench task types, metrics, input lengths and sample counts, and Table~\ref{tab:hyperparams} the LORE-KV configuration used unchanged throughout. LongBench, Needle-in-a-Haystack, and RULER results are analyzed in Sections~\ref{subsec:longbench}, \ref{subsec:niah}, and~\ref{subsec:ruler}, respectively; efficiency and mechanism analysis are reported in Sections~\ref{subsec:efficiency} and~\ref{subsec:mechanism}, and the ablations in Table~\ref{tab:ablation_summary} follow the protocol of Appendix~\ref{sec:ablation_components}. Full implementation details and LORE-KV hyperparameters are provided in Appendix~\ref{app:experimental_details}.

\textbf{The two reliability variants.} Every table reports both variants of Eq.~\eqref{eq:confidence} under an otherwise identical pipeline: \lorekvE, at $(\beta,\gamma,\lambda_D)=(1,1.3,0)$, whose numbers are complete, and \lorekvD, at $(1,1,0.3)$, whose end-to-end runs are in progress and are marked \tbd. The two share particle generation, the perturbation score, budget allocation and selection, so any difference between the rows comes from the reliability coefficients alone---but note that they differ in \emph{both} $\gamma$ and $\lambda_D$ (Proposition~1), so the pair is not a single-factor contrast; the single-factor cells $(1,0,0)$ and $(1,0,\lambda_D)$ are reported as ablations in Table~\ref{tab:ablation_summary}, and the two variants themselves are summarized in Appendix~\ref{app:reparam_grid}. Their score-level comparison on identical captured particles is already available and is reported in Appendix~\ref{app:divergence}.

\begin{table}[t]
\caption{LongBench results with Mistral-7B-Instruct-v0.2 across cache budgets. Baseline values are from \citet{geng2026accurate}; \tbd\ denotes runs in progress. LAQ++(8) is a compute-matched extension of Lookahead Q-Cache \citep{wang2025lookahead} run under our implementation at $B\in\{128,256,512\}$ (not evaluated at $B=1024$); it is reported for reference and excluded from the \best/\second\ highlighting.}
\label{tab:longbench_mistral}
\centering
\resizebox{\textwidth}{!}{%
\begin{tabular}{lccccccccccccccccc}
\toprule
& \multicolumn{3}{c}{Single-Doc. QA} 
& \multicolumn{3}{c}{Multi-Doc. QA} 
& \multicolumn{3}{c}{Summarization} 
& \multicolumn{3}{c}{Few-shot} 
& \multicolumn{2}{c}{Synthetic} 
& \multicolumn{2}{c}{Code} 
& \\
\cmidrule(lr){2-4}
\cmidrule(lr){5-7}
\cmidrule(lr){8-10}
\cmidrule(lr){11-13}
\cmidrule(lr){14-15}
\cmidrule(lr){16-17}
Method 
& NrtvQA & Qasper & MF-en 
& HotpotQA & 2WikiMQA & Musique 
& GovReport & QMSum & MultiNews 
& TREC & TriviaQA & SAMSum 
& PCount & PRe 
& Lcc & RB-P 
& Avg. \\
\midrule

Full Cache 
& 26.63 & 32.99 & 49.34 
& 42.77 & 27.35 & 18.77 
& 32.87 & 24.24 & 27.10 
& 71.00 & 86.23 & 42.96 
& 2.75 & 86.98 
& 55.33 & 52.87 
& 42.51 \\

% =====================================================================
% Budget = 128
% =====================================================================
\midrule
\multicolumn{18}{c}{\textit{Budget} $=128$} \\
\midrule

H2O 
& 21.19 & 21.66 & 38.60 
& 30.63 & 20.65 & 12.19 
& 20.65 & 22.42 & 21.81 
& 39.00 & 82.52 & 40.68 
& 2.98 & 79.56 
& 49.13 & 46.76 
& 34.40 \\

StreamingLLM 
& 16.61 & 14.74 & 31.40 
& 28.05 & 21.36 & 12.08 
& 18.44 & 18.91 & 19.26 
& 43.50 & 74.22 & 29.00 
& 2.75 & 31.65 
& 41.27 & 38.84 
& 27.63 \\

SnapKV 
& 19.17 & 21.40 & 42.93 
& 36.76 & 22.44 & 15.86 
& 19.16 & 21.84 & 21.55 
& 47.50 & 84.15 & 40.24 
& 2.30 & 68.26 
& 50.69 & 47.13 
& 35.09 \\

PyramidKV 
& 20.16 & 21.77 & 43.55 
& 36.78 & 23.12 & 14.39 
& 19.53 & 22.03 & 21.47 
& 51.00 & 84.62 & 40.24 
& 2.79 & 70.77 
& 50.57 & 46.53 
& 35.58 \\

Ada-KV 
& 21.79 & 23.03 & 47.07 
& 38.70 & 22.85 & 15.92 
& 19.94 & 23.05 & 21.84 
& 63.00 & 85.36 & 40.01 
& 3.20 & 72.13 
& 51.36 & 49.29 
& 37.41 \\

CriticalKV 
& 21.55 & 23.27 & 46.67 
& 38.95 & 24.95 & 16.85 
& 20.83 & 23.14 & 22.16 
& 65.00 & 85.53 & 40.02 
& \best{3.48} & 77.81 
& 53.58 & 50.98 
& 38.42 \\

AnDPro 
& \second{24.84} & 25.14 & 47.33 
& 39.79 & 23.99 & 17.34 
& 20.40 & \best{23.45} & 21.94 
& 67.50 & 85.59 & 40.62 
& 3.19 & 76.94 
& 53.77 & 51.55 
& 38.96 \\

LAQ++(8) 
& 25.62 & 27.21 & 46.16 
& 40.60 & 25.93 & 18.44 
& 21.60 & 23.07 & 22.42 
& 70.00 & 86.18 & 42.03 
& 3.51 & 74.81 
& 54.68 & 50.92 
& 39.57 \\

\textbf{\lorekvE\ (ours)} 
& 24.05 & \second{26.85} & \best{47.59} 
& \second{42.19} & \best{26.39} & \best{19.20} 
& \second{21.48} & 22.46 & \best{22.78} 
& \second{69.50} & \best{86.64} & \second{41.08} 
& \second{3.44} & \second{84.37} 
& \best{54.90} & \second{51.56} 
& \second{40.28} \\

\textbf{\lorekvD\ (ours)} 
& \best{25.14} & \best{27.94} & \second{47.34} 
& \best{42.26} & \second{25.91} & \second{18.71} 
& \best{21.54} & \second{23.35} & \second{22.61} 
& \best{70.00} & \second{86.33} & \best{41.24} 
& 3.03 & \best{84.51} 
& \second{54.24} & \best{52.03} 
& \best{40.39} \\

% =====================================================================
% Budget = 256
% =====================================================================
\midrule
\multicolumn{18}{c}{\textit{Budget} $=256$} \\
\midrule

H2O 
& 21.54 & 22.92 & 42.56 
& 31.07 & 22.53 & 13.76 
& 22.52 & 22.40 & 23.09 
& 40.50 & 84.20 & 40.77 
& \best{3.41} & 86.10 
& 50.98 & 48.17 
& 36.03 \\

StreamingLLM 
& 17.93 & 16.01 & 33.36 
& 30.71 & 21.30 & 10.08 
& 20.66 & 19.47 & 22.89 
& 53.50 & 73.59 & 29.22 
& 3.00 & 27.77 
& 42.30 & 39.87 
& 28.85 \\

SnapKV 
& 22.37 & 23.74 & 48.13 
& 38.56 & 22.43 & 15.66 
& 21.91 & 23.13 & 23.15 
& 61.50 & 85.45 & 41.42 
& 3.09 & 84.54 
& 53.22 & 50.24 
& 38.66 \\

PyramidKV 
& 20.09 & 24.00 & 47.33 
& 38.24 & 22.48 & 16.02 
& 21.40 & 22.45 & 22.63 
& 63.00 & 84.93 & 40.98 
& \second{3.40} & 82.48 
& 52.78 & 49.36 
& 38.22 \\

Ada-KV 
& 23.02 & 25.70 & \second{49.23} 
& 40.18 & 24.76 & 17.43 
& 22.17 & 23.21 & 23.48 
& 67.50 & 85.58 & 41.67 
& 2.99 & 84.65 
& 55.07 & 52.17 
& 39.93 \\

CriticalKV 
& 23.70 & 26.66 & 49.00 
& 40.05 & 25.12 & 17.86 
& 21.99 & 23.35 & 23.45 
& 68.00 & 86.19 & \second{42.07} 
& 2.92 & 85.81 
& 53.81 & 50.42 
& 40.03 \\

AnDPro 
& 25.71 & 28.17 & \best{49.31} 
& 40.04 & 25.81 & \second{18.68} 
& 21.66 & 23.64 & 23.63 
& \second{70.50} & \best{86.59} & \best{42.60} 
& 2.86 & 85.62 
& 55.74 & 52.87 
& 40.84 \\

LAQ++(8)
& 25.23 & 29.16 & 49.24 
& 41.70 & 26.85 & 18.62 
& 23.73 & 23.69 & 23.38 
& 70.50 & 86.24 & 42.54 
& 3.36 & 86.11 
& 55.59 & 52.49 
& 41.15 \\

\textbf{\lorekvE\ (ours)} 
& \second{25.88} & \second{29.58} & 48.10 
& \second{42.13} & \best{26.83} & \best{18.90} 
& \second{23.23} & \second{23.79} & \best{24.07} 
& \second{70.50} & \second{86.31} & 41.09 
& 3.27 & \second{89.09} 
& \best{56.15} & \second{53.12} 
& \second{41.38} \\

\textbf{\lorekvD\ (ours)} 
& \best{26.49} & \best{30.22} & 47.96 
& \best{42.17} & \second{26.56} & 18.63 
& \best{23.26} & \best{24.28} & \second{23.98} 
& \best{70.78} & 86.14 & 41.18 
& 3.04 & \best{89.17} 
& \second{55.78} & \best{53.38} 
& \best{41.44} \\

% =====================================================================
% Budget = 512
% =====================================================================
\midrule
\multicolumn{18}{c}{\textit{Budget} $=512$} \\
\midrule

H2O 
& 21.72 & 26.03 & 44.81 
& 32.33 & 23.16 & 14.86 
& 23.65 & 22.84 & 24.70 
& 42.00 & 85.22 & 41.57 
& \best{3.40} & 86.45 
& 53.04 & 49.68 
& 37.22 \\

StreamingLLM 
& 18.76 & 17.17 & 37.09 
& 30.21 & 21.64 & 9.93 
& 24.44 & 20.00 & \best{25.57} 
& 62.00 & 72.36 & 29.95 
& 2.48 & 18.17 
& 43.70 & 40.13 
& 29.60 \\

SnapKV 
& 24.60 & 27.81 & 48.98 
& 39.46 & 25.25 & 16.98 
& 23.70 & 22.96 & 24.37 
& 67.00 & 85.88 & 41.26 
& 2.78 & 86.56 
& 54.81 & 51.71 
& 40.26 \\

PyramidKV 
& 23.23 & 27.94 & 48.87 
& 40.50 & 24.36 & 16.74 
& 23.22 & 23.16 & 24.37 
& 67.00 & 85.73 & 41.74 
& \second{3.16} & 85.67 
& 54.16 & 50.34 
& 40.01 \\

Ada-KV 
& 23.88 & 29.05 & 48.79 
& 40.44 & 25.30 & 17.66 
& 23.47 & 23.47 & 24.48 
& 69.50 & 86.39 & 41.38 
& 2.94 & 87.56 
& 56.58 & 52.88 
& 40.86 \\

CriticalKV 
& 23.83 & 28.67 & \second{49.13} 
& 40.18 & 26.18 & 18.19 
& 23.52 & 23.57 & 24.28 
& 70.00 & 86.39 & \second{42.16} 
& 2.82 & 87.55 
& 56.51 & 52.64 
& 40.98 \\

AnDPro 
& 24.87 & \second{30.99} & \best{49.58} 
& 41.34 & 27.01 & 18.44 
& 23.40 & 23.76 & 25.01 
& \second{71.00} & \best{86.74} & \best{43.65} 
& 2.83 & 86.23 
& \best{57.46} & 53.61 
& \second{41.62} \\

LAQ++(8)
& 25.49 & 30.92 & 49.72 
& 41.50 & 26.84 & 19.20 
& 25.67 & 24.04 & 25.31 
& 71.00 & 86.43 & 43.14 
& 2.90 & 85.27 
& 56.80 & 53.54 
& 41.70 \\

\textbf{\lorekvE\ (ours)} 
& \second{25.58} & \best{31.14} & 48.98 
& \second{41.82} & \best{27.41} & \best{19.12} 
& \second{25.46} & \second{23.97} & \second{25.05} 
& \second{71.00} & \second{86.49} & 41.92 
& 2.95 & \second{88.20} 
& \second{56.85} & \second{53.82} 
& \best{41.86} \\

\textbf{\lorekvD\ (ours)} 
& \best{25.98} & 30.90 & 48.89 
& \best{41.85} & \second{27.23} & \second{18.94} 
& \best{25.48} & \best{24.30} & 24.99 
& \best{71.18} & 86.38 & 41.98 
& 2.80 & \best{88.25} 
& 56.61 & \best{53.99} 
& \best{41.86} \\

% =====================================================================
% Budget = 1024
% =====================================================================
\midrule
\multicolumn{18}{c}{\textit{Budget} $=1024$} \\
\midrule

H2O 
& 23.90 & 28.62 & 46.46 
& 37.03 & 24.74 & 15.04 
& 25.30 & 23.11 & 25.92 
& 46.00 & 85.93 & 41.80 
& \best{3.24} & 86.57 
& 54.46 & 51.01 
& 38.70 \\

StreamingLLM 
& 19.42 & 21.69 & 41.75 
& 32.40 & 22.18 & 11.18 
& \best{27.13} & 21.09 & \best{26.59} 
& 67.00 & 71.79 & 30.11 
& 2.88 & 16.57 
& 44.82 & 39.76 
& 31.02 \\

SnapKV 
& 25.47 & 29.57 & \second{49.33} 
& 40.90 & 25.53 & \best{19.01} 
& 25.94 & 23.89 & 26.21 
& 69.50 & 86.48 & 42.10 
& \second{2.98} & \best{88.56} 
& 55.57 & 51.92 
& 41.44 \\

PyramidKV 
& 24.21 & 29.86 & 48.93 
& 40.75 & 25.05 & 18.77 
& 25.73 & 24.06 & 25.65 
& 68.50 & 86.31 & 42.25 
& 2.97 & 87.17 
& 54.75 & 52.10 
& 41.07 \\

Ada-KV 
& 25.60 & 31.15 & 49.02 
& 41.47 & 27.07 & 18.77 
& 25.40 & 23.84 & 25.92 
& 70.50 & \second{86.55} & 43.07 
& 2.54 & 87.14 
& 57.27 & \best{53.85} 
& 41.82 \\

CriticalKV 
& \second{26.03} & 31.50 & 48.39 
& 41.73 & \second{27.10} & 18.84 
& 25.30 & 24.10 & 25.99 
& \second{71.00} & 86.30 & \second{43.16} 
& 2.66 & \second{87.23} 
& 56.69 & 53.20 
& 41.83 \\

AnDPro 
& 25.78 & \best{32.54} & \best{49.35} 
& 42.20 & 26.94 & 18.63 
& 25.33 & 24.01 & 26.29 
& \second{71.00} & 86.30 & \best{43.46} 
& 2.68 & 86.98 
& 57.18 & \second{53.77} 
& 42.03 \\

\textbf{\lorekvE\ (ours)} 
& 25.98 & \second{31.81} & 49.26 
& \second{43.96} & \best{27.17} & \second{18.91} 
& 26.98 & \second{24.12} & \second{26.41} 
& \second{71.00} & \best{86.60} & 42.30 
& 2.68 & 86.49 
& \best{57.83} & \second{53.77} 
& \best{42.20} \\

\textbf{\lorekvD\ (ours)} 
& \best{26.17} & 31.10 & 49.22 
& \best{43.97} & 27.08 & 18.82 
& \second{26.99} & \best{24.28} & 26.38 
& \best{71.09} & 86.54 & 42.33 
& 2.61 & 86.51 
& \second{57.71} & \best{53.85} 
& \second{42.17} \\

\bottomrule
\end{tabular}%
}
\end{table}

\subsection{Evaluation on LongBench}
\label{subsec:longbench}

Table~\ref{tab:longbench_mistral} reports results on LongBench with Mistral-7B-Instruct-v0.2; the corresponding Llama-3.1-8B-Instruct results are reported in Table~\ref{tab:longbench_llama}. \lorekvE\ achieves the best average performance at every cache budget. The gain is largest under aggressive compression: at $B=128$ it reaches $40.28$ versus $38.96$ for the strongest baseline, AnDPro, while the gap narrows to $42.21$ versus $42.03$ at $B=1024$. Both tables also report \textbf{LAQ++(8)}, a compute-matched extension of Lookahead Q-Cache \citep{wang2025lookahead} that scores tokens from $8$ local observation-window queries together with $8$ forward Q-cache entries, matching \lorekvE's lookahead length $T{=}8$; results are available at $B\in\{128,256,512\}$. Even against this equal-compute baseline, \lorekvE\ remains higher on average at every reported budget (e.g., $40.28$ vs.\ $39.57$ on Mistral and $46.11$ vs.\ $38.70$ on Llama at $B{=}128$, narrowing to $41.86$ vs.\ $41.70$ and $48.06$ vs.\ $40.39$ at $B{=}512$), though LAQ++(8) occasionally edges ahead on individual Mistral subtasks at $B{=}128$. Appendix~\ref{app:divergence} isolates the reliability term itself at score level. We do not yet attribute this margin between the ensemble estimator and the choice of reliability channel: every $\lambda_D>0$ cell is still \tbd, and the two reported configurations are $(\beta,\gamma,\lambda_D)=(1,1.3,0)$ and $(1,1,0.3)$, which differ in \emph{two} coefficients rather than one (Proposition~1), so they do not form a single-factor comparison.

We additionally measure how much cache a baseline requires to match the quality obtained by \lorekvE, reported in Table~\ref{tab:budget_equivalence}. At $B=128$, the strongest baseline requires approximately $208$ retained tokens to match \lorekvE, corresponding to a $38.5\%$ cache saving; the saving remains $31.1\%$ at $B=512$. This provides an operational interpretation of the accuracy gains: under tight budgets, improved token selection translates directly into substantially lower cache requirements. Figure~\ref{fig:budget_equivalence} shows the same trend as a continuous curve in behavioral fidelity, sustained across the whole budget sweep rather than at four discrete points.

\begin{table}[t]
\caption{Baseline cache budgets required to match \lorekvE\ on LongBench with Mistral-7B-Instruct-v0.2, with corresponding cache savings.}
\label{tab:budget_equivalence}
\centering
\begin{tabular}{lcccc}
\toprule
\lorekvE\ budget & 128 & 256 & 512 & 1024 \\
\midrule
Equivalent Ada-KV budget & $\sim$332 & $\sim$745 & $>$1024 & $>$1024 \\
Equivalent CriticalKV budget & $\sim$307 & $\sim$709 & $>$1024 & $>$1024 \\
Equivalent AnDPro budget & \best{$\sim$208} & \best{$\sim$414} & \best{$\sim$743} & $>$1024 \\
\midrule
Budget saving vs.\ best baseline (\%) & 38.5 & 38.1 & 31.1 & $>0$ \\
\bottomrule
\end{tabular}
\end{table}

\subsection{Evaluation on Needle-in-a-Haystack}\label{subsec:niah}
Needle-in-a-Haystack tests whether eviction preserves evidence needed by future decoding queries. We place the target at varying depths in $4$K--$32$K contexts and retain $B=1024$ tokens with Llama-3.1-8B-Instruct. As shown in the depth$\times$length grid of Figure~\ref{fig:niah-heatmap}, retrieval is near saturation, and the methods mainly differ at extreme depths, where prompt-end queries are least representative of the evidence later required. \lorekvE\ achieves the highest exact-match average ($0.977$ vs.\ $0.955$ for CriticalKV and $0.932$ for AnDPro), with the same ordering under word recall. The gain therefore comes from preventing a small number of complete retrieval failures in precisely the cases where the prompt-side query reference is weakest. The corresponding \lorekvD\ run is in progress (\tbd); because the two variants differ only in the scalar particle weights, they retain the same number of tokens and share the same latency at this budget, so any difference between them would be a retention-quality difference alone. Figure~\ref{fig:retained_token_composition} gives the token-level counterpart, showing that \lorekvE\ retains answer-bearing tokens more often than the prompt-side criteria at the same budget.

\begin{table}[t]
\caption{RULER accuracy averaged over 13 tasks at 4K, 8K, and 16K context lengths.
\textbf{Avg.} is the mean across lengths.}
\label{tab:ruler_main}
\centering
\begin{tabular}{lcccccccc}
\toprule
& \multicolumn{4}{c}{Mistral-7B-Instruct-v0.2}
& \multicolumn{4}{c}{Llama-3.1-8B-Instruct} \\
\cmidrule(lr){2-5}\cmidrule(lr){6-9}
Method & 4K & 8K & 16K & Avg. & 4K & 8K & 16K & Avg. \\
\midrule

\multicolumn{9}{c}{\textit{Budget} $=128$} \\
\midrule
CriticalKV
& 43.83 & 38.75 & 36.08 & 39.55
& 50.72 & 48.85 & 45.06 & 48.21 \\

AnDPro
& 54.85 & 50.06 & 45.20 & 50.04
& 60.06 & 58.42 & 52.80 & 57.09 \\

\textbf{\lorekvE\ (ours)}
& \best{58.40} & \best{55.77} & \best{51.05} & \best{55.07}
& \best{64.27} & \best{63.43} & \best{56.05} & \best{61.25} \\

\textbf{\lorekvD\ (ours)}
& \second{58.31} & \second{55.74} & \second{51.00} & \second{55.02}
& \second{64.20} & \second{63.32} & \second{55.98} & \second{61.17} \\

\midrule
\multicolumn{9}{c}{\textit{Budget} $=256$} \\
\midrule
CriticalKV
& 63.89 & 56.74 & 51.95 & 57.53
& 65.55 & 60.74 & 58.25 & 61.51 \\

AnDPro
& \second{71.05} & 64.08 & 57.91 & 64.35
& 71.62 & 66.98 & 63.25 & 67.28 \\

\textbf{\lorekvE\ (ours)}
& \best{71.51} & \best{65.52} & \best{59.94} & \best{65.66}
& \best{72.22} & \best{68.34} & \best{65.80} & \best{68.78} \\

\textbf{\lorekvD\ (ours)}
& 69.77 & \second{64.56} & \second{58.98} & \second{64.44}
& \second{72.10} & \second{68.22} & \second{65.73} & \second{68.69} \\

\midrule
\multicolumn{9}{c}{\textit{Budget} $=512$} \\
\midrule
CriticalKV
& 79.26 & 67.75 & 63.65 & 70.22
& 79.25 & 69.02 & 67.38 & 71.88 \\

AnDPro
& \second{83.42} & 72.06 & \second{67.12} & \second{74.20}
& 82.71 & 72.58 & 70.28 & 75.19 \\

\textbf{\lorekvE\ (ours)}
& \best{83.69} & \best{72.82} & \best{68.28} & \best{74.93}
& \best{83.12} & \best{73.46} & \best{71.60} & \best{76.06} \\

\textbf{\lorekvD\ (ours)}
& 81.67 & \second{72.18} & 64.98 & 72.94
& \second{83.05} & \second{73.35} & \second{71.53} & \second{75.98} \\

\midrule
\multicolumn{9}{c}{\textit{Budget} $=1024$} \\
\midrule
CriticalKV
& 88.78 & 83.05 & 80.25 & 84.03
& 88.35 & 83.26 & \best{81.72} & 84.44 \\

AnDPro
& \second{90.97} & \second{85.31} & \second{82.06} & \second{86.11}
& 93.38 & 85.94 & 79.37 & 86.23 \\

\textbf{\lorekvE\ (ours)}
& \best{91.11} & \best{85.74} & \best{82.75} & \best{86.53}
& \best{93.60} & \best{86.40} & \second{79.98} & \best{86.66} \\

\textbf{\lorekvD\ (ours)}
& 89.88 & 84.51 & 81.53 & 85.31
& \second{93.50} & \second{86.29} & 79.87 & \second{86.55} \\

\bottomrule
\end{tabular}
\end{table}

\subsection{Evaluation on RULER: Sweeping Context Length at a Fixed Task} \label{subsec:ruler}
RULER provides a controlled test of the query-reference hypothesis by instantiating the same $13$ tasks at $4$K, $8$K, and $16$K contexts while keeping the task and lookahead cost fixed. As shown in Table~\ref{tab:ruler_main}, \lorekvE\ achieves the best overall average in all eight model--budget settings and the best per-length result in $23$ of $24$ comparisons; the per-task breakdowns are Tables~\ref{tab:ruler_mistral_full} and~\ref{tab:ruler_llama_full}, which locate the gain in the aggregation-style tasks. The gains are largest where the prompt-side reference is expected to be weakest: under tight budgets, the margin over AnDPro is $5.05$ points on Mistral and $3.52$ on Llama at $B=128$, but falls to about $0.4$ points at $B=1024$; with increasing context length, the gain also generally grows, e.g., from $3.58$ points at $4$K to $5.85$ at $16$K on Mistral with $B=128$. Since the lookahead cost is unchanged across these settings, these trends support a retention-criterion effect rather than an extra-compute effect: future-query-aware scoring matters most when the cache is highly constrained and the fixed prompt-end window becomes less representative of later decoding queries. The only per-length exception is Llama at $B=1024$, $16$K, where CriticalKV scores $81.72$ versus $79.98$ for \lorekvE, consistent with the smaller differences between eviction criteria once the cache budget is relatively generous. The \lorekvD\ rows of Table~\ref{tab:ruler_main} are \tbd; since RULER separates the criteria most sharply at $B{=}128$ and hardly at all at $B{=}1024$, that is where the two variants can be told apart, and it is the budget at which the \lorekvE-versus-\lorekvD\ comparison of Appendix~\ref{app:reparam_grid} is run.

As a further compute-matched reference, \textbf{LAQ++(8)} \citep{wang2025lookahead} reaches RULER averages of only $51.57$, $62.81$, and $69.83$ on Llama-3.1-8B-Instruct at $B=128,256,512$, respectively, versus $61.25$, $68.78$, and $76.06$ for \lorekvE\ (Table~\ref{tab:ruler_main})---a wider margin than on LongBench. Mistral-7B-v0.2 results, the $B=1024$ setting, and the per-length and per-task breakdowns are not available for this baseline, so it is reported here in text rather than added to Table~\ref{tab:ruler_main} or Tables~\ref{tab:ruler_mistral_full}--\ref{tab:ruler_llama_full}.

\subsection{Memory and Latency} \label{subsec:efficiency}
At a fixed cache budget $B$, all methods have the same steady-state decoding memory and per-token cost. LORE-KV adds only a one-time lookahead overhead: the $M$ particles are decoded as a single batched continuation of length $T$ using the prefetched prompt KV cache. This cost is paid once before final generation and is independent of output length; the reliability term adds no decoding of its own (Appendix~\ref{app:overhead_details}), so both variants have the same cost profile. The full procedure is given in Algorithm~\ref{alg:lorekv}, while Appendix~\ref{app:overhead_details} details the computational accounting and compute-matched baseline control.

Because \textbf{LAQ++(8)} \citep{wang2025lookahead} is evaluated with the same lookahead compute as \lorekvE\ ($T{=}8$), it isolates the effect of \emph{how} that compute is used rather than how much is spent. \lorekvE\ still leads on average at every reported budget and dataset against this equal-compute baseline (LongBench: $40.28$ vs.\ $39.57$ on Mistral and $46.11$ vs.\ $38.70$ on Llama at $B{=}128$, narrowing but persisting to $41.86$ vs.\ $41.70$ and $48.06$ vs.\ $40.39$ at $B{=}512$; RULER on Llama: $61.25$ vs.\ $51.57$ at $B{=}128$ and $76.06$ vs.\ $69.83$ at $B{=}512$; Section~\ref{subsec:longbench} and Section~\ref{subsec:ruler}). The advantage is therefore consistent and robust to matching test-time compute, indicating that it comes from the multi-trajectory, reliability-weighted estimator rather than from spending more test-time compute.

\begin{figure}[t]
  \centering
  \includegraphics[width=\textwidth]{figures_combined/fig3_cache_distribution.pdf}
  \caption{\textbf{Attention-output fidelity.} Cosine similarity between retained-cache and full-cache attention outputs on future queries, shown as distributions and layer-wise means. Mistral-7B-Instruct-v0.2, $B=128$.}
  \label{fig:cache_distribution}
\end{figure}

\begin{figure}[t]
  \centering
  \includegraphics[width=\textwidth]{figures_combined/fig6_mismatch_vs_gain.pdf}
  \caption{\textbf{Gain versus prompt-to-future query mismatch.} Per-example gains in oracle retention recall and output fidelity as mismatch increases. Mistral-7B-Instruct-v0.2, $B=128$, $M=4$.}
  \label{fig:mismatch_vs_gain}
\end{figure}

\subsection{Mechanism Analysis: Locating the Gain} \label{subsec:mechanism}
We evaluate three predictions: P1 (\emph{operating point}), P2 (\emph{saturation}), and P3 (\emph{locality}); Appendix~\ref{app:datasets} restates them in full and Appendix~\ref{app:mechanism_protocol} details the protocol used to measure them below. Sections~\ref{subsec:longbench}--\ref{subsec:ruler} establish P1; here we test P2 and P3 by directly measuring the query-reference mismatch identified in Section~\ref{sec:introduction}. Let $\mathcal{Q}_{\mathrm{future}}$ denote the queries that the compressed cache actually serves (Appendix~\ref{app:kv_background}). Since these queries are unavailable at eviction time, we record them afterwards from full-cache generation and use them only as ground truth for analysis. All results use Mistral-7B-Instruct-v0.2 at $B=128$; the full protocol is provided in Appendix~\ref{app:mechanism_protocol}, with additional analyses in Appendix~\ref{app:experiments_extra}. Figure~\ref{fig:future_query_alignment} shows that response-side lookahead queries better approximate $\mathcal{Q}_{\mathrm{future}}$ than the prompt-end window: on LongBench, coverage improves from $0.788$ to $0.826$ with one rollout and saturates at $0.828$-$0.829$ with additional rollouts. This supports P2: the main gain comes from replacing input-side with response-side queries, as in Section~\ref{subsec:lookahead}, rather than simply increasing rollout compute. The same saturation appears in Table~\ref{tab:ablation_summary}. Figure~\ref{fig:cache_distribution} further shows that the resulting cache better preserves the full-cache output targeted by Eq.~\eqref{eq:eviction_objective} and Eq.~\eqref{eq:projected_score}, increasing mean cosine similarity from $0.873$ to $0.905$ on LongBench and from $0.929$ to $0.944$ on RULER. Figure~\ref{fig:layer_head_analysis} shows that this improvement extends across the network, supporting the projected deletion criterion of Section~\ref{subsec:deletion}.

To test P3, we measure each example's prompt-to-future query mismatch,
$\operatorname{Mis}(X)=1-\operatorname{Cov}(\mathcal{Q}_{\mathrm{future}};\{\mathcal{Q}^{l,h}_{\mathcal{W}}\})$,
and relate it to the \lorekvE\ gain in Figure~\ref{fig:mismatch_vs_gain}. Larger mismatch consistently corresponds to larger gains: Spearman correlations are $\rho=0.28$ and $0.57$ for oracle retention recall, and $\rho=0.53$ and $0.51$ for output fidelity on LongBench and RULER, respectively. Together with the saturation in Figure~\ref{fig:future_query_alignment} (P2) and the budget- and length-dependent results in Section~\ref{subsec:ruler} (P1), this supports the predicted mechanism: gains arise from correcting the query reference rather than merely adding rollout compute. Figure~\ref{fig:future_query_alignment} additionally separates the two reliability variants: the entropy- and disagreement-weighted estimators are drawn as separate curves over the same nested particle prefixes, so the vertical distance between them is the entire effect of the penalty, and it is small next to the step from prompt-side to response-side queries. Additional checks are reported in Appendix~\ref{app:experiments_extra}, including Figures~\ref{fig:score_comparison}, \ref{fig:confidence_calibration}, \ref{fig:particle_divergence}, \ref{fig:retained_token_composition}, \ref{fig:layer_head_analysis}, and~\ref{fig:budget_equivalence}, with reliability weighting defined in Eq.~\eqref{eq:confidence}.

\section{Related Work}
\label{sec:related_work}
\subsection{Prompt-side KV eviction} \label{subsec:rw_eviction}
Most prior methods estimate token importance from prefill or prompt-side statistics. StreamingLLM preserves attention sinks and recent tokens \citep{xiao2023streamingllm}; H2O and Scissorhands use accumulated attention \citep{zhang2023h2o,liu2023scissorhands}; TOVA uses current-query attention \citep{oren2024transformers}; FastGen applies head-specific retention patterns \citep{ge2024fastgen}; and SnapKV uses a prompt-end observation window \citep{li2024snapkv}. PyramidKV and CAKE improve layer-wise budget allocation \citep{cai2024pyramidkv,qin2025cake}, while Ada-KV, HeadKV, DynamicKV, SAGE-KV, LAVa, and LU-KV refine allocation across heads, layers, or token groups \citep{feng2024adakv,fu2024headkv,zhou2024dynamickv,wang2025sagekv,shen2025lava,tang2026lukv}. DefensiveKV further introduces robust aggregation against poorly estimated queries \citep{feng2025taming}. These methods improve retention or budget allocation but still estimate importance using prompt-side queries rather than the future queries that consume the compressed cache.

\subsection{Future-query estimation and output-aware scoring} \label{subsec:rw_future}
Several methods improve the query reference. Lookahead Q-Cache uses queries from a single surrogate response \citep{wang2025lookahead}; Expected Attention analytically estimates future attention \citep{devoto2025expected}; and DapQ constructs pseudo-queries from future positional information \citep{tian2026dapq}. LookaheadKV, KVP, globally calibrated retention methods, and IndexMem instead learn future token utility or compact memory representations \citep{ahn2026lookaheadkv,moschella2026learning,bui2026make,yang2026indexmem}. KVzip, FastKVzip, and ContrastKV target query-agnostic or multi-query robustness \citep{kim2026fastkvzip,kim2026kvzip,chen2026contrastkv}. A complementary line improves the importance score itself. CriticalKV incorporates value information \citep{guo2024criticalkv}; AnDPro uses projected attention-output information \citep{geng2026accurate}; and LaProx models projected values and inter-head effects \citep{mai2026laprox}. KeyDiff and $L_2$-norm methods use key geometry \citep{park2026keydiff,devoto2024simple}, while LagKV, GraphKV, R-KV, and MomentKV exploit temporal, graph, redundancy, or directional structure \citep{liang2025lagkv,li2025graphkv,cai2025rkv,li2026momentkv}. Recent analysis also shows that value-aware corrections depend strongly on the quality of the query reference \citep{zhang2026valueaware}. Future-aware methods improve the query reference but rely on single, analytic, or learned estimates; output-aware methods improve deletion scoring but remain prompt-side. Our reliability weighting connects this line to uncertainty estimation for free-form generation, where sequence likelihood, predictive entropy, and agreement across sampled continuations are the standard signals \citep{kuhn2023semantic,farquhar2024detecting}; the difference here is that agreement is measured directly between the sampled \emph{query} trajectories, i.e.\ in the representation the eviction score consumes, rather than between the generated strings.

\subsection{Orthogonal KV-cache compression} \label{subsec:rw_orthogonal}
Other approaches reduce KV-cache cost along dimensions orthogonal to token eviction. Quantization methods such as KIVI, KVQuant, QAQ, MiKV, MiniKV, RotateKV, KVTuner, MixKVQ, PM-KVQ, and LeanKV reduce precision or combine sparsity with mixed precision \citep{liu2024kivi,hooper2024kvquant,dong2024qaq,yang2024notokenleft,sharma2024minikv,su2025rotatekv,li2025kvtuner,zhang2025mixkvq,liu2025pmkvq,10.1145/3731569.3764810}. GEAR, MiniCache, KVMerger, xKV, inter-layer methods, and transform coding exploit low-rank, cross-layer, mergeable, or coded representations \citep{kang2024gear,liu2024minicache,wang2024kvmerger,chang2025xkv,ma2024interlayer,staniszewski2026kv}. KVServe targets communication-efficient serving \citep{liu2026kvserve}, while DMS, ThinKV, RaaS, LazyEviction, R-KV, RocketKV, LongFlow, and CONF-KV address cache growth during long decoding \citep{lancucki2025hyperscaling,ramachandran2025thinkv,hu2025raas,zhang2025lazyevict,cai2025rkv,behnam2025rocketkv,su2026longflow,li2026confkv}. These methods are complementary because they change cache precision, structure, serving, or decode-time maintenance rather than the prompt-time token-importance criterion. Aggressive eviction can also remove instructions or answer-bearing evidence \citep{yang2024notokenleft,chen2025pitfalls,garcia2026protection}. Overall, prior methods either approximate future queries incompletely or improve deletion scoring while remaining prompt-side; the key gap is to combine multi-trajectory future-query estimation with future-aware deletion scoring.

\section{Conclusion} \label{sec:conclusion}

LORE-KV reframes KV-cache eviction from identifying which prompt tokens were important to estimating which tokens will remain useful to future decoding queries, and it does so by \emph{sampling} that future rather than predicting it: a small Monte-Carlo ensemble drawn from the frozen model replaces the prompt-side query reference. This addresses a central limitation of prior attention- and output-aware methods: their scores can be unreliable when computed from an unrepresentative prompt-side query reference. Making the ensemble an estimator also makes its weighting a design choice with a stated target. Reaching that target required separating a confound: because the particles are drawn from the model itself, a sampled token's log-probability already contains the entropy of the distribution it came from, so likelihood and uncertainty cannot be priced independently in the obvious parametrization. Centering the per-token score restores that independence and gives three coefficients---branch atypicality, context uncertainty, and ensemble-relative trajectory disagreement---which we compared under an otherwise identical pipeline. Our results show that the dominant effect is the substitution of response-side for prompt-side queries: lookahead is most useful under tight budgets and long contexts, provides little benefit when prompt-end queries already approximate future queries well, and the weighting rule is a second-order refinement on top of it. This motivates adaptively enabling lookahead only under high query mismatch, otherwise using prompt-only eviction, with future extensions to decode-time cache updates, adaptive particle counts driven by observed ensemble disagreement, and orthogonal compression methods.


\subsection*{AI Use Statement}

We used generative AI only to improve the grammar, clarity, and concision of author-written text, and coding assistants for routine implementation tasks such as plotting, table generation, and code refactoring. AI was not used to develop the research ideas, design the method or experiments, or generate and interpret results. All experimental results were obtained from actual runs or cited prior work. The authors reviewed and verified all AI-assisted content and take full responsibility for the final paper.

\subsection*{Ethics Statement}

This work uses no human subjects, personally identifying data, or new data collection; experiments rely on publicly released LongBench, RULER, and Needle-in-a-Haystack benchmarks and the public Mistral-7B-Instruct-v0.2 and Llama-3.1-8B-Instruct checkpoints under their respective licenses. LORE-KV reduces KV-cache memory and bandwidth without adding new generative capabilities. However, aggressive eviction may remove instructions, boundaries, or evidence needed by later queries, potentially degrading instruction following or factual grounding; Appendix~\ref{app:composition} analyzes this behavior and Appendix~\ref{app:limitations} documents cases where the method is neutral or harmful. Deployments under tight cache budgets should therefore be validated on safety-relevant application prompts rather than average benchmark accuracy alone.

\subsection*{Reproducibility Statement}

LORE-KV is training-free and modifies no model weights. The method is fully specified by Eqs.~\eqref{eq:lookahead_cont}--\eqref{eq:gather} in Section~\ref{sec:methodology}, with pseudocode in Algorithm~\ref{alg:lorekv}; all hyperparameters and default switches are listed in Table~\ref{tab:hyperparams}, and the two reliability variants are the points $(\beta,\gamma,\lambda_D)=(1,1.3,0)$ and $(1,1,0.3)$ of Eq.~\eqref{eq:confidence}, so either is reproduced from the other by changing those coefficients. A run of the earlier $(\lambda_H,\eta)$ parametrization is reproduced exactly by $(\beta,\gamma,\lambda_D)=(\eta,\eta(1+\lambda_H),\eta\lambda_D)$ (Proposition~1). Appendix~\ref{app:implementation} documents grouped-query-attention head mapping and scoring over $\mathcal{Q}^{l,h}_{\mathrm{score}}$. Evaluation details include LongBench datasets, metrics, and sample counts in Table~\ref{tab:longbench_datasets}, RULER construction and scoring in Appendix~\ref{app:ruler_results}, and the budget grid, chunk size, and observation window in Section~\ref{subsec:setup}, which also distinguishes our runs from values taken from prior work. All executed LongBench configurations are reported in Table~\ref{tab:ablation_summary}, with RULER ablations in Table~\ref{tab:ruler_ablation}; Appendix~\ref{app:mechanism_protocol} specifies the mechanism-analysis capture procedure and plotted quantities, with sample sizes and lookahead budgets summarized in Table~\ref{tab:figure_recipes}.

\bibliography{main}
\bibliographystyle{iclr2027_conference}

\appendix

\section{Background: The KV Cache and the Eviction Problem} \label{app:kv_background}

\textbf{KV caching and memory cost.}
A decoder-only Transformer \citep{vaswani2017attention} generates tokens autoregressively, projecting each hidden state into a query $\mathbf{q}_t$, key $\mathbf{k}_t$, and value $\mathbf{v}_t$, with attention weights determined by $\mathbf{q}_t^\top\mathbf{k}_j/\sqrt{d}$. Since past keys and values remain fixed during inference, they are stored in the \emph{KV cache}, enabling a parallel \emph{prefill} over the prompt followed by token-by-token decoding. For $L$ layers, $H$ KV heads, head dimension $d$, sequence length $n$, batch size $b$, and $p$ bytes per element, the cache requires $\mathrm{Mem}_{\mathrm{KV}} = 2bLHdnp$ bytes, where the factor $2$ accounts for keys and values.Thus, memory grows linearly with context length even at batch size one. For example, consider a 70B-class model with $L=80$ layers, $H=8$ KV heads, $d=128$, a $128$K-token context ($n=131{,}072$), $b=1$, and $p=2$ bytes (BF16). Its KV cache alone requires $\mathrm{Mem}_{\mathrm{KV}}\approx42.9$ GB. For comparison, a 70B-parameter model stored in BF16 requires roughly $140$ GB for its weights, so a single 128K context adds KV-cache memory equal to about $31\%$ of the model's weight memory.

\textbf{The eviction problem.}
\label{app:kv_eviction_problem}
Given a cache budget $B\ll n$, eviction retains index sets
$\mathcal{R}^{l,h}\subseteq\{1,\ldots,n\}$ with
$|\mathcal{R}^{l,h}|\le B_{l,h}$ and removes the remaining KV pairs. Ideally, the retained cache should preserve the behavior of the full-cache model during future decoding:
\begin{equation}
\min_{\{\mathcal{R}^{l,h}\}}
\mathbb{E}_{\mathbf{q}\sim\mathcal{Q}_{\mathrm{future}}}
\left[
\mathrm{dist}\!\left(
\mathrm{output}(\mathbf{q};\text{full cache}),
\mathrm{output}(\mathbf{q};\text{retained cache})
\right)
\right].
\label{eq:eviction_objective}
\end{equation}
Future decoding queries $\mathcal{Q}_{\mathrm{future}}$ are unknown at eviction time, so methods must estimate which cached tokens will matter. Prior work uses prompt-side or predicted queries with attention- or output-based scores, whereas LORE-KV uses confidence-weighted lookahead queries and ranks tokens by their leave-one-out output impact. Section~\ref{subsec:mechanism} evaluates this approximation against the actual future queries observed after generation.

\section{Method Details} \label{app:method_details}
This appendix discusses the detailed design, pseudocode, cost accounting, and implementation of the methodology in Section~\ref{sec:methodology}. 

\subsection{Implementation: Protected Set, Grouped-Query Attention, Fused Kernels}
\label{app:implementation}
\label{app:setup_details}
\label{app:gqa}
\label{app:kernels}

\textbf{Protected positions.}
The optional set $\mathcal{B}$ holds the positions of prompt-structure markers: chat-template delimiters (\verb!<|system|>!, \verb!<|user|>!, \verb!<|assistant|>!, \verb![INST]!, \verb![/INST]!), tool-call markers and document separators. These are read off the tokenized prompt and always retained; if the template exposes none, $\mathcal{B}=\emptyset$. Following prior work \citep{xiao2023streamingllm,li2024snapkv}, this keeps instructions and prompt boundaries from being removed under aggressive eviction. 

\textbf{Grouped-query attention.}
Under GQA/MQA, multiple query heads share each KV head. We therefore score each cached pair $(\mathbf{k}^{l,h}_{i},\mathbf{v}^{l,h}_{i})$ using all query heads in its group $\mathcal{G}(h)$ and aggregate their deletion effects from Eq.~\eqref{eq:deletion_effect}. Throughout, $H$ denotes the number of KV heads, and cache budgets are assigned at the KV-head level.

\textbf{Fused attention kernels.}
Kernels such as FlashAttention never materialize the $n\times n$ attention matrix a perturbation score needs. We avoid kernel surgery by computing scores only for the scoring-query set $\mathcal{Q}^{l,h}_{\mathrm{score}}$ of Eq.~\eqref{eq:score_queries}, whose size $|\mathcal{Q}^{l,h}_{\mathrm{score}}|=w+MT$ is far below $n$. For those few queries we form the attention weights against the cached keys explicitly, leaving the fused kernel untouched for the rest of prefill and for final decoding. Eviction scoring is thus a small side computation that never requires full-matrix attention on the main inference path.

\subsection{Relation to Prior Surrogate-Query and Output-Aware Methods} \label{app:lookahead_details}
LORE-KV inherits two components and adds one. Response-side surrogate queries for estimating future cache usage come from Lookahead Q-Cache \citep{wang2025lookahead,ahn2026lookaheadkv}, and output-aware deletion scoring beyond attention alone comes from AnDPro and its relatives \citep{guo2024criticalkv,geng2026accurate,mai2026laprox}. What is added is that the surrogate is a \emph{Monte-Carlo ensemble} rather than a single continuation, together with the reliability weighting that turns the ensemble into an estimator with a stated target distribution (Section~\ref{subsec:confidence}). Concretely, Lookahead Q-Cache is the $M{=}1$ special case of the query construction and AnDPro is the $\lambda_L{=}0$ special case of the score, so both are reachable from Eqs.~\eqref{eq:score_queries} and~\eqref{eq:final_score} by switch settings, and both are priced in Tables~\ref{tab:ablation_summary} and~\ref{tab:ablation_summary}. We weight the trajectories using model likelihood together with either predictive entropy \citep{kuhn2023semantic,farquhar2024detecting}, which is \lorekvE, or ensemble disagreement measured in query space, which is \lorekvD. The two inherited components are common to both variants, so the $M{=}1$ and $\lambda_L{=}0$ special cases are shared: at $M{=}1$ the softmax of Eq.~\eqref{eq:particle_weight} normalizes a single scalar away and $\hat D_m$ is identically zero, so \lorekvE\ and \lorekvD\ coincide exactly with Lookahead Q-Cache-style scoring, and the variants can only differ at $M\ge3$. Figure~\ref{fig:particle_divergence} measures both signals against the same per-particle oracle utility and ranks them against each other on identical particles, while Figure~\ref{fig:confidence_calibration} shows that either confidence is only modestly correlated with that utility, motivating soft weighting rather than single-trajectory selection.

\begin{algorithm}[t]
\caption{LORE-KV eviction procedure.}
\label{alg:lorekv}
\begin{algorithmic}[1]

\Require frozen LLM $f_\theta$; prompt $X=(x_1,\ldots,x_n)$; window $w$;
particles $M$, length $T$; layer budgets $\{B_l\}_{l=1}^{L}$;
core constants $\epsilon,\beta,\gamma,\lambda_D,\lambda_L$;
plug-in switches $\alpha,\kappa,\tau,\lambda_k,\lambda_v$
\Comment{all switches $=0$ gives the core method;
$(\beta,\gamma,\lambda_D)=(1,1.3,0)$ is \lorekvE\ and $(1,1,0.3)$ is \lorekvD}

\Ensure compressed KV cache $(\mathbf{K}^\star,\mathbf{V}^\star)$

\State $\big(\{\mathbf{K}^{l,h},\mathbf{V}^{l,h}\},
\{\mathcal{Q}^{l,h}_{\mathcal{W}}\}\big)
\gets \operatorname{Prefill}(f_\theta,X)$
\Comment{build the KV cache and collect prompt-end queries}

\State $\mathcal{P}\gets\{1\}\cup\mathcal{W}\cup\mathcal{B}$;\quad
$\mathcal{C}\gets\{1,\ldots,n\}\setminus\mathcal{P}$
\Comment{separate protected and candidate tokens}

\Statex \textcolor{blue}{\textit{\% Stage 1: estimate future decoding queries}}

\For{$m=1,\ldots,M$ \textbf{in parallel}}

    \State $Y^{(m)}\sim q(\cdot\mid X)$,
    $\ |Y^{(m)}|=T$;\ store
    $\mathcal{Q}^{(m),l,h}_{\mathrm{look}}\ \forall(l,h)$,\ $\{c^{(m)}_u,H^{(m)}_u\}_{u=1}^{T}$
    \Comment{sample a lookahead; store its queries and per-step $c_u$ (Eq.~\eqref{eq:centered_score}), $H_u$}

\EndFor

\State $\mathbf{D}\gets\textsc{Disagree}
\big(\{\mathcal{Q}^{(m),l,h}_{\mathrm{look}}\}_{m=1}^{M}\big)$
by Eq.~\eqref{eq:disagreement};\ \ $\hat{\mathbf{D}}\gets\operatorname{zscore}(\mathbf{D})$ if $M\ge3$ else $\mathbf{0}$
\Comment{disagreement from the stored queries; no extra decoding}

\State $\bar c_m\gets\frac1T\sum_{u}c^{(m)}_u$,\ \
$\bar H_m\gets\frac1T\sum_{u}H^{(m)}_u$,\ \
$R_m\gets\beta\bar c_m-\gamma\bar H_m-\lambda_D\hat D_m$
\ \ $\forall m$
\Comment{means, not sequence totals (Eq.~\eqref{eq:confidence})}

\State $\boldsymbol{\omega}\gets
\operatorname{softmax}(\mathbf{R})\in\Delta^{M-1}$,
$\ \mathbf{R}=(R_m)_{m=1}^{M}$
\Comment{convert reliability scores into trajectory weights}

\State $\mathcal{Q}^{l,h}_{\mathrm{score}}
\gets
\mathcal{Q}^{l,h}_{\mathcal{W}}
\cup
\bigcup_{m=1}^{M}\mathcal{Q}^{(m),l,h}_{\mathrm{look}}$
\Comment{combine prompt and lookahead queries}

\Statex \textcolor{blue}{\textit{\% Stage 2: score the impact of removing each token}}

\For{layer $l=1,\ldots,L$ \textbf{and} head $h=1,\ldots,H$}

    \For{candidate $i\in\mathcal{C}$ \textbf{and}
    query $\mathbf{q}\in\mathcal{Q}^{l,h}_{\mathrm{score}}$}

        \State $(a^{l,h}_{q,i},\mathbf{o}^{l,h}_{q})
        \gets
        \operatorname{Attn}
        (\mathbf{q},\mathbf{K}^{l,h},\mathbf{V}^{l,h})$
        \Comment{compute attention weight and output}

        \State $\Delta\mathbf{o}^{l,h}_{q,i}
        \gets
        \dfrac{a^{l,h}_{q,i}}
        {1-a^{l,h}_{q,i}+\epsilon}
        \big(\mathbf{v}^{l,h}_{i}-\mathbf{o}^{l,h}_{q}\big)$

        \State $e^{l,h}_{q,i}
        \gets
        \left\|
        \mathbf{W}^{O}_{l,h}
        \Delta\mathbf{o}^{l,h}_{q,i}
        \right\|_2^2$
        \Comment{measure the effect of removing token $i$}

    \EndFor

    \State $S^{l,h}_{i}
    \gets
    \rho(\mathcal{E}^{l,h}_{i,\mathcal{W}})
    +
    \lambda_L
    \sum_{m=1}^{M}
    \omega_m\,
    \rho(\mathcal{E}^{l,h}_{i,m})
    \quad \forall i\in\mathcal{C}$
    \Comment{combine prompt and weighted lookahead importance}

    \Statex \textcolor{blue}{\textit{\% Stage 3: allocate cache budget and retain tokens}}

    \State $\{B_{l,h}\}_{h=1}^{H}
    \gets
    \operatorname{Alloc}
    \big(B_l,\{E_{l,h}\}_{h=1}^{H};\tau\big)$
    \Comment{divide the layer budget across heads}

    \For{head $h=1,\ldots,H$}

        \State $\mathcal{R}^{l,h}
        \gets
        \operatorname{Select}
        (\mathbf{S}^{l,h},B_{l,h})$
        \Comment{keep the highest-scoring tokens}

        \State $\mathcal{R}^{l,h}_{\mathrm{final}}
        \gets
        \mathcal{P}\cup\mathcal{R}^{l,h}$

        \State $(\mathbf{K}^{l,h}_{\star},
        \mathbf{V}^{l,h}_{\star})
        \gets
        \operatorname{Gather}
        (\mathbf{K}^{l,h},
        \mathbf{V}^{l,h},
        \mathcal{R}^{l,h}_{\mathrm{final}})$
        \Comment{build the compressed cache}

    \EndFor

\EndFor

\State \textbf{discard} lookahead states
\State \Return $(\mathbf{K}^{\star},\mathbf{V}^{\star})$
\Comment{use the compressed cache for final decoding}

\end{algorithmic}
\end{algorithm}


\begin{table*}[t]
\centering
\caption{Average inference time per sample (seconds) across datasets, backbones, and cache budgets. Lower is better.}
\label{tab:runtime_comparison}
\small
\setlength{\tabcolsep}{5.0pt}
\renewcommand{\arraystretch}{1.08}
\begin{tabular}{lllcccc}
\toprule
\textbf{Backbone} & \textbf{Dataset} & \textbf{Method}
& \textbf{128} & \textbf{256} & \textbf{512} & \textbf{1024} \\
\midrule

\multirow{9}{*}{Mistral}
& \multirow{3}{*}{LongBench}
& AndPro      & 7.8799 & 7.4885 & 7.4355 & 7.5338 \\
&& CriticalKV & 6.6867 & 6.3991 & 6.3982 & 6.5276 \\
&& \lorekvE\ & 17.5342 & 19.5496 & 19.6767 & 17.7930 \\
\cmidrule(lr){2-7}

& \multirow{3}{*}{RULER}
& AndPro      & 6.2360 & 4.5775 & 5.9707 & 3.1746 \\
&& CriticalKV & 6.1907 & 4.1849 & 4.9899 & 2.4039 \\
&& \lorekvE\ & 13.6877 & 11.1720 & 14.5537 & 7.2315 \\
\cmidrule(lr){2-7}

& \multirow{3}{*}{NIAH}
& AndPro      & 2.9842 & 2.8559 & 2.8779 & 2.9169 \\
&& CriticalKV & 2.3514 & 2.1441 & 2.1109 & 2.1961 \\
&& \lorekvE\ & 7.3873 & 7.4716 & 7.4072 & 7.3407 \\
\midrule

\multirow{9}{*}{LLaMA}
& \multirow{3}{*}{LongBench}
& AndPro      & 8.3347 & 8.2447 & 8.1548 & 8.0648 \\
&& CriticalKV & 7.0726 & 7.0443 & 7.0160 & 6.9877 \\
&& \lorekvE\ & 12.1815 & 14.4038 & 15.8882 & 13.1342 \\
\cmidrule(lr){2-7}

& \multirow{3}{*}{RULER}
& AndPro      & 4.3309 & 4.0226 & 3.7143 & 3.4060 \\
&& CriticalKV & 4.2994 & 3.6776 & 3.1042 & 2.5791 \\
&& \lorekvE\ & 6.6410 & 6.8587 & 6.3249 & 5.4202 \\
\cmidrule(lr){2-7}

& \multirow{3}{*}{NIAH}
& AndPro      & 2.6046 & 2.5137 & 2.5269 & 2.6172 \\
&& CriticalKV & 1.9267 & 1.9222 & 1.9450 & 1.8561 \\
&& \lorekvE\ & 7.0298 & 6.9553 & 7.0076 & 6.9816 \\
\bottomrule
\end{tabular}
\end{table*}


\subsection{Reliability Weighting: Derivations, Calibration, and Edge Cases}
\label{app:reliability_details}

This section collects the derivations and implementation details omitted from Section~\ref{subsec:confidence}.

\textbf{Centering and reparameterization.}
The centered score of Eq.~\eqref{eq:centered_score} removes the context-entropy baseline from the likelihood channel in conditional expectation. For $y_u\sim q(\cdot\mid h_u)$,
\begin{equation}
\mathbb{E}[c_u\mid h_u]
=\sum_y q(y\mid h_u)\log q(y\mid h_u)+H_u
=-H_u+H_u=0,
\label{eq:centering}
\end{equation}
so $\mathbb{E}[\bar c_m]=0$. Moreover, $S_t=\sum_{u\le t}c_u$ is a martingale with respect to the natural filtration. This centering should not be read as statistical independence between $\bar c_m$ and $\bar H_m$; it makes branch atypicality zero-mean conditional on the visited context so that uncertainty can be parameterized explicitly through $\gamma$.

\textbf{Proposition 1 (nesting).}
Because $\tfrac1T\sum_u\log q(y^{(m)}_u\mid h^{(m)}_u)=\bar c_m-\bar H_m$, the earlier score
$\tfrac1T\sum_u\log q-\lambda_H\bar H_m-\lambda_D\hat D_m$
is exactly
$\bar c_m-(1+\lambda_H)\bar H_m-\lambda_D\hat D_m$.
Thus uniform weighting is $(\beta,\gamma,\lambda_D)=(0,0,0)$; the earlier ``log-probability only'' setting is $(1,1,0)$, not a no-entropy control; \lorekvE\ is $(1,1.3,0)$; and \lorekvD\ is $(1,1,0.3)$. This is also why \lorekvE\ and \lorekvD\ are variant settings rather than a single-factor ablation.

\textbf{Variance and scale.}
For $u<v$, $c_u$ is $h_v$-measurable and $\mathbb{E}[c_uc_v]=\mathbb{E}[c_u\,\mathbb{E}[c_v\mid h_v]]=0$. Hence
\begin{equation}
\operatorname{Var}(\bar c_m)
=\frac{1}{T^2}\sum_{u=1}^{T}\mathbb{E}[V_u],
\qquad
V_u=\operatorname{Var}_{y\sim q(\cdot\mid h_u)}[\log q(y\mid h_u)],
\label{eq:varentropy}
\end{equation}
where $V_u$ is the varentropy. The scale of $\bar c_m$ can therefore vary across prompts. The optional normalization
$\tilde c_m=\bar c_m/\sqrt{\bar V_m/T}$, with $\bar V_m=\tfrac1T\sum_uV^{(m)}_u$, gives unit scale by construction; it is evaluated only as an ablation and is not part of the default method.

\textbf{Query-space disagreement.}
Predictive entropy is defined in token-probability space, whereas the eviction score consumes query directions. We therefore define the disagreement of particle $m$ as the mean cosine distance to the other particles at matched decoding positions, averaged over lookahead steps, layers, and query heads:
\begin{equation}
D_m=\frac{1}{M-1}\sum_{r\neq m}
\frac{1}{TLH}\sum_{u=1}^{T}\sum_{l=1}^{L}\sum_{h=1}^{H}
\left[1-\cos\!\left(\mathbf q^{(m),l,h}_u,\mathbf q^{(r),l,h}_u\right)\right].
\label{eq:disagreement}
\end{equation}
Queries compared at step $u$ occupy the same absolute position $n+u$ and therefore share the same rotary phase. We standardize disagreement across particles,
\begin{equation}
\hat D_m=\frac{D_m-\bar D}{\operatorname{sd}(D)},
\qquad
\bar D=\frac1M\sum_{r=1}^{M}D_r.
\label{eq:disagreement_z}
\end{equation}
For $M=1$, disagreement is undefined; for $M=2$, $D_1=D_2$ by symmetry and the standard deviation is zero. We therefore set $\hat D_m\equiv0$ for $M\le2$, so disagreement can affect weights only for $M\ge3$. For ragged trajectories, $\bar c_m$ and $\bar H_m$ use each particle's available steps, whereas $D_m$ is computed on the common prefix required for step-wise alignment (Appendix~\ref{app:datasets}). The disagreement computation reuses the stored lookahead queries and requires no additional decoding (Appendix~\ref{app:overhead_details}).

The entropy and disagreement coefficients are not directly comparable numerically because $\bar H_m$ is measured in nats while $\hat D_m$ is standardized. We therefore calibrate $\lambda_D$ using the realized across-particle spread of the two penalties rather than equating their raw coefficients.

\textbf{Softmax parameterization and estimator interpretation.}
No separate particle-weight temperature is introduced: multiplying Eq.~\eqref{eq:confidence} by a concentration parameter $\eta$ is equivalent to rescaling $(\beta,\gamma,\lambda_D)$, so $\eta$ is not separately identifiable.

With uniform weights, the lookahead term is the ordinary Monte-Carlo estimate of expected deletion effect under continuations sampled from $p_\theta(\cdot\mid X)$. When $\lambda_D=0$, each weight depends only on its own trajectory and the normalized weighted estimate has target
\[
\pi(Y)\propto p_\theta(Y\mid X)\exp\!\left(\beta\bar c(Y)-\gamma\bar H(Y)\right).
\]
To first order,
\[
\mathbb{E}_{\pi}[f]
=\mathbb{E}_{p_\theta}[f]
+\beta\operatorname{Cov}_{p_\theta}(f,\bar c)
-\gamma\operatorname{Cov}_{p_\theta}(f,\bar H)
+O(\|(\beta,\gamma)\|^2).
\]
Because $\mathbb{E}_{p_\theta}[\bar c]=0$, centering removes a zeroth-order shift from branch atypicality; it does not imply an unbiased finite-$\beta$ estimator or zero higher-order covariance. When $\lambda_D>0$, $\hat D_m$ depends on the full particle set, so the weights form an interacting-particle estimator rather than finite-$M$ importance weights for a fixed target law. A fixed disagreement functional is recovered only asymptotically as $M$ grows.

\textbf{Failure modes and safeguards.}
\lorekvE\ can assign substantial weight to a confidently wrong continuation because low predictive entropy does not guarantee correctness. \lorekvD\ can suppress such a trajectory when its query path is atypical relative to the sampled ensemble, but it does not detect a collectively wrong consensus and is unreliable when too few particles are available. Step-wise cosine disagreement can also overstate differences between semantically similar continuations with different tokenizations or lengths. These limitations motivate soft weighting, the bounded contribution of the lookahead term through $\lambda_L$, and the use of \lorekvE\ as the default until the disagreement variant is supported by complete end-to-end evidence. Figures~\ref{fig:confidence_calibration} and~\ref{fig:particle_divergence} report the empirical relation between these signals and oracle particle utility.

\subsection{Computational Overhead} \label{app:overhead_details}

At a fixed cache budget $B$, all methods retain the same number of KV pairs and therefore have matched steady-state cache memory and final-decoding cost; LORE-KV trades additional one-time wall-clock computation for more reliable token retention under the same memory budget. Specifically, the $M$ conditionally independent lookahead particles are generated as a single batched continuation of length $T$ from the shared prompt cache, and scoring is restricted to the $w+MT\ll n$ prompt-window and lookahead queries. The reliability term costs nothing on top of this: the entropy of \lorekvE\ is a by-product of the sampling logits, and the disagreement of \lorekvD\ is one $M\times M$ cosine reduction per layer over query tensors that are already materialized, i.e.\ $O(M^2LHTd)$ arithmetic with no additional memory and no additional decoding, which is negligible beside the $MT$ decode steps that produced those tensors. The two variants therefore have the same wall-clock profile, and Table~\ref{tab:runtime_comparison} applies to both. \textcolor{red}{Table~\ref{tab:runtime_comparison} quantifies this accuracy-latency trade-off: \lorekvE\ requires approximately $1.46\times$ to $2.77\times$ the per-sample wall-clock time of AnDPro across the reported settings, while preserving the same $B$-token cache for final decoding; the same numbers apply to \lorekvD\ by the cost argument above.} This additional pre-generation computation provides a more representative future-query reference and improves retention quality, particularly under aggressive compression where preserving task-relevant evidence is most difficult. Importantly, the overhead does not increase monotonically with $B$, consistent with its primary source being the fixed $(M,T)$ lookahead and scoring procedure rather than increased cache memory. Several controls further separate the improved retention criterion from additional computation alone: \textbf{(i)} all methods use the same $B$, matching steady-state memory and per-token final-decoding cost; \textbf{(ii)} despite fixed lookahead compute, the RULER gain over AnDPro decreases from $5.05$ to $0.42$ points on Mistral and from $3.52$ to $0.43$ on LLaMA as $B$ increases from $128$ to $1024$ (Table~\ref{tab:ruler_main}), query coverage saturates at small $M$ (Figure~\ref{fig:future_query_alignment}), and per-example gains increase with prompt-to-future query mismatch (Figure~\ref{fig:mismatch_vs_gain}); \textbf{(iii)} $\lambda_L=0$ (equivalently $M=0$) removes lookahead and recovers the prompt-window criterion (Table~\ref{tab:ablation_summary}); and \textbf{(iv)} a compute-matched control supplying the same lookahead continuations to Lookahead Q-Cache-style scoring directly equalizes test-time compute: the end-to-end result is \textbf{LAQ++(8)} \citep{wang2025lookahead}, reported in Tables~\ref{tab:longbench_mistral}--\ref{tab:longbench_llama} and Section~\ref{subsec:ruler}, where \lorekvE\ remains ahead on average despite the matched compute; the corresponding score-level analysis is reported in Appendix~\ref{app:score_shift}.

\subsection{Hyperparameters and Switches}
\label{app:hyperparams}
\label{app:score_details}

The core method exposes seven quantities: the observation window $w$, the lookahead count $M$ and length $T$, the three reliability coefficients $\beta$ (typicality), $\gamma$ (entropy) and $\lambda_D$ (trajectory disagreement), and the mixing coefficient $\lambda_L$; the constant $\epsilon$ only guards the renormalization in Eq.~\eqref{eq:deletion_effect}. The plug-ins add $\alpha,\kappa$ (robust aggregation), $\tau$ (adaptive head budgets) and $\lambda_k,\lambda_v$ (diversity selection); we write the robust-aggregation top-fraction as $\kappa$ because $\beta$ now names the typicality coefficient of Eq.~\eqref{eq:confidence}. Each plug-in recovers its default exactly at zero: $\lambda_L{=}0\Leftrightarrow M{=}0$ (no lookahead), $\alpha{=}0$ (mean aggregation), $\tau{=}0$ (uniform budgets), $\lambda_k{=}\lambda_v{=}0$ (top-score selection), which is what makes the block-level ablation protocol well-posed as detailed in Table~\ref{tab:hyperparams}. The reliability block is different in kind: $(\beta,\gamma,\lambda_D)$ is a \emph{plane} that all three regimes---\lorekvE\ at $(1,1.3,0)$, \lorekvD\ at $(1,1,0.3)$, and log-probability only at $(1,1,0)$---occupy as distinct points (Proposition~1), rather than a set of switches, so a run of any of them reproduces the others by changing coefficients, with no re-decoding when the particles are already captured. Note that the no-weighting control is $(0,0,0)$, and the no-\emph{entropy} control is $(1,0,0)$; neither coincides with the earlier ``log-probability only'' setting, which is $(1,1,0)$.

There is no particle-weight temperature $\eta$: it is not identifiable alongside $(\beta,\gamma,\lambda_D)$ (Section~\ref{subsec:confidence}) and has been absorbed into them.
\begin{table}[t]
\caption{LORE-KV hyperparameters and default settings.}
\label{tab:hyperparams}
\centering
\begin{tabular}{llc}
\toprule
Symbol & Meaning & Default \\
\midrule
$w$ & prompt-end observation-window size & 32 \\
$M$ & number of lookahead particles & 4 \\
$T$ & lookahead particle length & 8 \\
$\beta$ & typicality coefficient in Eq.~\eqref{eq:confidence} & $1.0$ \\
$\gamma$ & entropy coefficient in Eq.~\eqref{eq:confidence} & $1.3$ (\lorekvE) / $1$ (\lorekvD) \\
$\lambda_D$ & disagreement coefficient in Eq.~\eqref{eq:confidence} & $0$ (\lorekvE) / $0.3$ (\lorekvD) \\
$\lambda_L$ & lookahead mixing coefficient in Eq.~\eqref{eq:final_score} & 1.0 \\
$\epsilon$ & renormalization guard in Eq.~\eqref{eq:deletion_effect} & $10^{-6}$ \\
\midrule
$\kappa,\alpha$ & robust aggregation (Appendix~\ref{sec:ablation_components}) & $1,\,0$ \\
$\tau$ & adaptive head budgets (Appendix~\ref{sec:ablation_components}) & $0$ \\
$\lambda_k,\lambda_v$ & diversity selection (Appendix~\ref{sec:ablation_components}) & $0$ \\
\bottomrule
\end{tabular}
\end{table}

\section{Plug-In Components and the Ablation Protocol}
\label{sec:ablation_components}

The core method deliberately keeps all downstream operations simple, using mean aggregation, uniform head-wise budgets, and top-score selection. Building on these defaults, we consider three optional plug-ins motivated by established designs in prior work. Each is controlled by a separate switch, allowing its incremental effect to be isolated relative to the fixed core configuration. The plug-ins act on $\rho$, $\operatorname{Alloc}$ and $\operatorname{Select}$ and are therefore orthogonal to the reliability coefficients: the executed rows sit on the \lorekvE\ core, and their \lorekvD\ counterparts---the same switch settings at $(\beta,\gamma,\lambda_D)=(1,1,0.3)$---are the \tbd\ block of Table~\ref{tab:ablation_summary}. Results are reported in Table~\ref{tab:ablation_summary} for LongBench, Table~\ref{tab:ruler_ablation} for RULER, and Table~\ref{tab:ablation_summary} for the complete configuration grid.

\textbf{A. Robust aggregation ($\alpha,\kappa$).}
\label{subsec:robust_aggregation}
A plain mean can suppress a token that is important for only a small subset of future queries, motivating a robust aggregation strategy in the spirit of DefensiveKV \citep{feng2025taming}. We therefore replace the mean with a top-$\kappa$ average, $\operatorname{TopAvg}_\kappa(\mathcal{E})=\frac{1}{K_\kappa}\sum_{e\in\operatorname{TopK}(\mathcal{E},K_\kappa)} e$, where $K_\kappa=\lceil \kappa|\mathcal{E}|\rceil$, and interpolate it with the standard mean as $\rho(\mathcal{E})=\alpha\operatorname{TopAvg}_\kappa(\mathcal{E})+(1-\alpha)\operatorname{Mean}(\mathcal{E})$, with $\alpha\in[0,1]$. (This top-fraction was written $\beta$ in an earlier draft; it is renamed $\kappa$ here because $\beta$ now denotes the typicality coefficient of Eq.~\eqref{eq:confidence}.) Substituting $\rho(\mathcal{E})$ for the mean in Eq.~\eqref{eq:final_score} yields the robust variant; $\alpha=0$ exactly recovers the core configuration, while larger $\alpha$ increases sensitivity to high-importance query subsets and provides an additional safeguard against a misleading particle.

\textbf{B. Adaptive head-wise budgets ($\tau$).}
\label{subsec:adaptive_budget}
The core method splits the layer budget evenly across heads. Adaptive allocation is inherited from Ada-KV \citep{feng2024adakv} and the layer-wise scheme of PyramidKV \citep{cai2024pyramidkv}; our one variation is that the allocation mass is each head's total future-query perturbation,
\begin{equation}
E_{l,h}=\sum_{i\in\mathcal{C}}S^{l,h}_{i},
\label{eq:pert_mass}
\end{equation}
which is what $\operatorname{Alloc}$ of Eq.~\eqref{eq:alloc_select} consumes. Its $\tau>0$ branch is
\begin{equation}
B_{l,h}=\left\lfloor
B_l\frac{(E_{l,h}+\epsilon)^{\tau}}{\sum_{h'=1}^{H}(E_{l,h'}+\epsilon)^{\tau}}
\right\rfloor,
\label{eq:adaptive_budget}
\end{equation}
with $\tau\ge0$ setting how aggressively capacity shifts toward high-impact heads; $\tau=0$ returns $\lfloor B_l/H\rfloor$. This separates ``better scores'' from ``better placement of a fixed budget.'' Figure~\ref{fig:layer_head_analysis} supplies the premise: both $E_{l,h}$ and the benefit of the future-query term are distributed very unevenly across depth and heads.

\textbf{C. Diversity-aware selection ($\lambda_k,\lambda_v$).}
\label{subsec:diversity}
Top-$B_{l,h}$ selection can in principle retain many near-duplicate keys or values. As an ablation we make $\operatorname{Select}$ a greedy rule inspired by redundancy-aware selection such as GraphKV \citep{li2025graphkv}: starting from $\mathcal{R}^{l,h}=\emptyset$, each newly selected token is
\begin{equation}
i^\star=\arg\max_{i\in\mathcal{C}\setminus\mathcal{R}^{l,h}}
\left[
S^{l,h}_{i}
-\lambda_k\max_{j\in\mathcal{R}^{l,h}}\cos(\mathbf{k}^{l,h}_{i},\mathbf{k}^{l,h}_{j})
-\lambda_v\max_{j\in\mathcal{R}^{l,h}}\cos(\mathbf{v}^{l,h}_{i},\mathbf{v}^{l,h}_{j})
\right],
\label{eq:diversity_select}
\end{equation}
repeated until $|\mathcal{R}^{l,h}|=B_{l,h}$; $\lambda_k=\lambda_v=0$ recovers plain top-$B_{l,h}$. The known failure mode is that a cosine penalty can discard evidence that is repeated but genuinely important, as when several supporting passages mention the same entity. The penalty acts on the perturbation score rather than on raw key/value features, so a duplicate that still drives a large future-query perturbation survives; and $\lambda_k,\lambda_v$ are small, so redundancy only breaks ties between tokens of comparable importance. The output-perturbation score may already suppress redundancy on its own, since a duplicated token whose deletion is compensated by its near-copy has a small $\Delta\mathbf{o}^{l,h}_{q,i}$ by construction.

\textbf{Ablations of the core itself.}
\label{subsec:core_ablations}
Three axes stress the claimed contribution directly. \textbf{(i)~Lookahead budget:} $M\in\{0,1,2,4,8\}$ at fixed $T{=}8$, where $M{=}0$ (equivalently $\lambda_L{=}0$) collapses to a prompt-window output-perturbation method; saturation at small $M$ is prediction P2. Varying $T$ is a planned axis we have not yet run. \textbf{(ii)~Reliability weighting:} the $(\beta,\gamma,\lambda_D)$ cells of Table~\ref{tab:ablation_summary}, which price each channel of Eq.~\eqref{eq:confidence} separately. At $\lambda_D{=}0$ these are the missing uniform control $(0,0,0)$, the no-entropy control $(1,0,0)$, the stronger-entropy probe $(1,2,0)$ and the two published cells $(1,1,0)$ and $(1,1.3,0)$; at $\lambda_D{>}0$ they add $(1,1,0.3)$ and the single-factor cell $(1,0,0.3)$ that puts disagreement in place of entropy rather than on top of it. The two cells the method itself ships as variants---\lorekvE\ and \lorekvD---are Appendix~\ref{app:reparam_grid}; the rest are controls. All cells run at uniform head budgets and top-$K$ selection, everything else held fixed, and re-weight nested prefixes of one $8$-particle draw per prompt so no cell differs by test-time compute. \textbf{(iii)~Perturbation objective:} we replace Eq.~\eqref{eq:projected_score} with accumulated attention mass, the unprojected perturbation norm $\|\Delta\mathbf{o}^{l,h}_{q,i}\|_2^2$, under the identical lookahead queries, providing the compute-matched control described in Appendix~\ref{app:overhead_details}.

\section{Experimental Setup and Additional Results}
\label{app:experimental_details}

\subsection{Protocol, Datasets and Baselines}\label{app:datasets}
We evaluate LORE-KV on LongBench, RULER, and Needle-in-a-Haystack while testing three mechanism-level predictions: gains should be larger under tighter budgets and longer contexts \textbf{(P1)}, saturate after a few lookahead trajectories \textbf{(P2)}, and increase with prompt window future-query mismatch \textbf{(P3)}. LongBench covers 16 English tasks across QA, summarization, few-shot learning, synthetic retrieval, and code, with input lengths from $1{,}235$ to $18{,}409$ tokens; Table~\ref{tab:longbench_datasets} reports task metrics, average lengths, and sample counts. RULER provides a controlled complement by evaluating 13 fixed-length tasks, including eight needle variants, variable tracking, common- and frequent-word extraction, and two QA tasks; because its question appears at the prompt end, the observation window is a comparatively strong proxy for future queries. Needle-in-a-Haystack varies target depth and context length to test whether answer-bearing evidence survives eviction.

\textbf{Baselines and sourcing.}
Table~\ref{tab:longbench_mistral} compares StreamingLLM \citep{xiao2023streamingllm}, H2O \citep{zhang2023h2o}, SnapKV \citep{li2024snapkv}, PyramidKV \citep{cai2024pyramidkv}, Ada-KV \citep{feng2024adakv}, CriticalKV \citep{guo2024criticalkv}, and AnDPro \citep{geng2026accurate}, with Full Cache as the uncompressed reference. StreamingLLM through Ada-KV use prompt-side attention statistics, whereas CriticalKV and AnDPro incorporate value- or output-aware information while remaining prompt-side. LongBench results with Mistral follow the published protocol of \citet{geng2026accurate}; CriticalKV, AnDPro, and \lorekvE\ are re-run under one implementation for RULER, Needle-in-a-Haystack, and LongBench with Llama; \lorekvD\ runs in that same implementation and is pending at submission time. Lookahead Q-Cache \citep{wang2025lookahead} is additionally included as the compute-matched end-to-end baseline \textbf{LAQ++(8)} in Tables~\ref{tab:longbench_mistral} and~\ref{tab:longbench_llama} (Section~\ref{subsec:longbench}) and, where available, in Section~\ref{subsec:ruler}; the corresponding score-level compute-matched control remains a separate diagnostic, analyzed in Appendix~\ref{app:score_shift}.

\textbf{Implementation.}
Our implementation builds on the Ada-KV framework and uses a single NVIDIA H100 80\,GB GPU. All methods use the same retained cache size at a fixed budget, and adjacent-token grouping follows SnapKV and Ada-KV. Unless stated otherwise, LORE-KV uses $M{=}4$, $T{=}8$, $\lambda_L{=}1.0$, and disables all plugins
($\alpha{=}\tau{=}\lambda_k{=}\lambda_v{=}0$), thereby isolating the score in Eq.~\eqref{eq:final_score}. The reliability coefficients are the only difference between the two variants: $(\beta,\gamma,\lambda_D){=}(1,1.3,0)$ for \lorekvE\ and $(1,1,0.3)$ for \lorekvD.

Lookahead particles are sampled at temperature $1$ with no top-$k$ or top-$p$ truncation, so the sampling law is $p_\theta$ itself and the centering of Eq.~\eqref{eq:centered_score} is exact. The implementation computes $c_u$, $H_u$ and the varentropy $V_u$ from the \emph{same} shaped logits the sampler consumed, so a future change of sampling law cannot silently break $\mathbb{E}[c_u\mid h_u]=0$; entropy and varentropy use masked reductions, since outside a truncated support $q\log q$ would otherwise be $0\cdot(-\infty)$. Per-step statistics are stored, not only trajectory means, at a cost of $MT$ floats per prompt. Disagreement is computed over every decoder layer (a layer stride is available but unused), from the lookahead query buffers the scoring stage already holds, and must be evaluated before the cache is compressed because compression releases those buffers. If a particle terminates early, the $D_m$ comparison is truncated to the shortest trajectory in the ensemble so that the step-wise pairing in Eq.~\eqref{eq:disagreement} stays valid; $\bar c_m$ and $\bar H_m$ need no such rule, since each averages over its own particle's steps.

\begin{table}[t]
\caption{LongBench English datasets, metrics, average input lengths, and sample counts \citep{bai2024longbench}.}
\label{tab:longbench_datasets}
\centering
\begin{tabular}{lllrlr}
\toprule
Task & Task type & Metric & Avg.\ len & Language & \#Samples \\
\midrule
NarrativeQA & Single-Doc. QA & F1 & 18{,}409 & EN & 200 \\
Qasper & Single-Doc. QA & F1 & 3{,}619 & EN & 200 \\
MultiFieldQA-en & Single-Doc. QA & F1 & 4{,}559 & EN & 150 \\
HotpotQA & Multi-Doc. QA & F1 & 9{,}151 & EN & 200 \\
2WikiMultihopQA & Multi-Doc. QA & F1 & 4{,}887 & EN & 200 \\
MuSiQue & Multi-Doc. QA & F1 & 11{,}214 & EN & 200 \\
GovReport & Summarization & Rouge-L & 8{,}734 & EN & 200 \\
QMSum & Summarization & Rouge-L & 10{,}614 & EN & 200 \\
MultiNews & Summarization & Rouge-L & 2{,}113 & EN & 200 \\
TREC & Few-shot Learning & Accuracy & 5{,}177 & EN & 200 \\
TriviaQA & Few-shot Learning & F1 & 8{,}209 & EN & 200 \\
SAMSum & Few-shot Learning & Rouge-L & 6{,}258 & EN & 200 \\
PassageCount & Synthetic & Accuracy & 11{,}141 & EN & 200 \\
PassageRetrieval-en & Synthetic & Accuracy & 9{,}289 & EN & 200 \\
LCC & Code & Edit Sim & 1{,}235 & Python/C\#/Java & 500 \\
RepoBench-P & Code & Edit Sim & 4{,}206 & Python/Java & 500 \\
\bottomrule
\end{tabular}
\end{table}


\begin{table}[t]
\caption{LongBench results with Llama-3.1-8B-Instruct across cache budgets. CriticalKV, AnDPro, and LORE-KV are matched runs; other baseline values are from \citet{geng2026accurate}. \tbd\ denotes runs in progress. LAQ++(8) is a compute-matched extension of Lookahead Q-Cache \citep{wang2025lookahead} run under our implementation at $B\in\{128,256,512\}$ (not evaluated at $B=1024$); it is reported for reference and excluded from the \best/\second\ highlighting.}
\label{tab:longbench_llama}
\centering
\resizebox{\textwidth}{!}{%
\begin{tabular}{lccccccccccccccccc}
\toprule
Method & NrtvQA & Qasper & MF-en & HotpotQA & 2WikiMQA & Musique & GovReport & QMSum & MultiNews & TREC & TriviaQA & SAMSum & PCount & PRe & Lcc & RB-P & Avg. \\
\midrule
Full Cache & 30.12 & 46.60 & 56.41 & 58.10 & 49.01 & 32.52 & 34.04 & 25.37 & 27.06 & 73.00 & 91.90 & 43.46 & 7.06 & 100.00 & 62.13 & 51.80 & 49.29 \\
\midrule
\multicolumn{18}{c}{\textit{Budget} $=128$} \\
\midrule
\multicolumn{18}{c}{\textit{Budget} $=128$} \\
\midrule
SnapKV & 25.54 & 31.40 & 50.99 & 55.17 & 45.50 & 27.88 & 20.77 & 23.41 & 21.05 & 47.50 & 90.46 & 40.76 & \second{8.00} & \second{99.50} & 56.88 & 46.50 & 43.21 \\
PyramidKV & \second{27.83} & 32.98 & 51.45 & \second{56.72} & 44.67 & \second{30.54} & 21.89 & 22.90 & 21.67 & 63.00 & 90.45 & 40.06 & \second{8.00} & \best{100.00} & 55.37 & 43.79 & 44.46 \\
Ada-KV & 25.06 & 34.50 & 52.27 & 56.29 & \second{47.57} & 29.06 & 21.17 & 23.36 & 21.51 & 62.00 & \best{92.17} & 40.84 & \second{8.00} & \second{99.50} & 58.45 & 48.72 & 45.03 \\
CriticalKV & 26.00 & 34.96 & 51.79 & 56.01 & 46.53 & 28.95 & 21.31 & 23.14 & 21.33 & 62.16 & 89.51 & 40.74 & 7.59 & 97.89 & 57.85 & 48.28 & 44.63 \\
AnDPro & 25.25 & 35.85 & \best{52.85} & 54.35 & 44.41 & 27.51 & 22.82 & \best{23.60} & 22.20 & \second{69.00} & 92.00 & 41.50 & 7.02 & \best{100.00} & \best{62.26} & \best{56.34} & 46.06 \\
LAQ++(8) \citep{wang2025lookahead} & 28.65 & 10.65 & 26.04 & 24.23 & 21.56 & 15.67 & 23.50 & 23.74 & 22.75 & 72.00 & 91.95 & 42.19 & 8.37 & 94.81 & 62.18 & 50.95 & 38.70 \\
\textbf{\lorekvE\ (ours)} & 27.38 & \second{37.66} & 51.14 & \best{57.37} & \best{48.47} & \best{30.55} & \second{23.77} & \best{23.60} & \best{23.14} & 68.04 & 90.46 & \second{41.77} & 7.34 & 97.90 & 58.99 & 50.15 & \second{46.11} \\
\textbf{\lorekvD\ (ours)} & \best{29.00} & \best{39.20} & \second{52.83} & 54.76 & 45.31 & 30.20 & \best{23.86} & \second{23.48} & \second{22.26} & \best{71.00} & \second{92.09} & \best{42.75} & \best{8.01} & \best{100.00} & \second{61.63} & \second{55.76} & \best{47.01} \\

\midrule
\multicolumn{18}{c}{\textit{Budget} $=256$} \\
\midrule
SnapKV & 26.92 & 38.19 & 52.62 & 57.03 & 46.75 & 29.63 & 23.07 & 23.92 & 22.82 & 58.50 & 91.83 & 40.81 & 7.75 & 99.50 & 59.95 & 48.63 & 45.50 \\
PyramidKV & \second{28.46} & 39.36 & 54.30 & 56.77 & 45.50 & 30.74 & 23.84 & 23.88 & 22.91 & 69.00 & 91.09 & 40.74 & \best{7.88} & 99.50 & 56.94 & 45.18 & 46.01 \\
Ada-KV & 27.04 & 39.24 & \second{54.41} & 57.36 & \second{47.53} & 31.35 & 23.06 & 23.81 & 23.43 & 69.00 & \best{92.50} & \second{41.23} & 7.67 & \best{100.00} & 61.45 & 49.52 & 46.79 \\
CriticalKV & 25.45 & 39.60 & 52.90 & 55.72 & 46.05 & 29.97 & 23.10 & 23.77 & 22.77 & 68.28 & 90.48 & 41.08 & 7.48 & 97.55 & 60.19 & 48.25 & 45.79 \\
AnDPro & 27.23 & 42.01 & \best{54.62} & \best{57.87} & 47.43 & 31.92 & 23.53 & 24.22 & 23.54 & \second{71.87} & \second{92.37} & \best{41.47} & 7.66 & \second{99.81} & \second{62.45} & \second{51.67} & 47.48 \\
LAQ++(8) \citep{wang2025lookahead} & 30.25 & 12.43 & 26.63 & 25.77 & 22.83 & 18.45 & 25.07 & 23.67 & 23.75 & 72.00 & 91.97 & 43.17 & 6.93 & 94.59 & 64.07 & 53.38 & 39.69 \\
\textbf{\lorekvE\ (ours)} & 28.26 & \second{44.01} & 53.43 & \second{57.73} & \best{48.78} & \best{32.30} & \second{25.63} & \best{24.93} & \best{24.10} & 71.43 & 91.03 & 39.98 & 7.39 & 98.52 & 61.70 & 51.28 & \second{47.53} \\
\textbf{\lorekvD\ (ours)} & \best{29.16} & \best{44.86} & 54.37 & 56.29 & 47.03 & \second{32.11} & \best{25.68} & \second{24.86} & \second{23.61} & \best{73.07} & 91.93 & 40.52 & \second{7.76} & 99.68 & \best{63.16} & \best{54.38} & \best{48.03} \\

\midrule
\multicolumn{18}{c}{\textit{Budget} $=512$} \\
\midrule
SnapKV & 28.40 & 41.52 & 55.40 & 57.66 & 47.48 & 30.95 & 24.92 & 24.46 & 24.54 & 68.50 & \second{92.33} & 41.98 & 7.33 & 99.50 & 61.25 & 50.39 & 47.29 \\
PyramidKV & 29.15 & 43.04 & \second{55.55} & 57.71 & 47.67 & 31.27 & 26.05 & 24.08 & 24.54 & 70.50 & \best{92.67} & 41.52 & 7.33 & \best{100.00} & 59.96 & 47.51 & 47.41 \\
Ada-KV & 26.92 & 42.68 & 55.20 & 57.47 & 48.06 & 31.27 & 24.95 & 24.17 & 24.40 & 71.00 & 92.14 & \best{42.35} & 7.33 & \best{100.00} & \second{62.26} & 50.62 & 47.55 \\
CriticalKV & 26.00 & 42.53 & 53.90 & 56.62 & 46.76 & 31.13 & 24.26 & 23.76 & 23.81 & 70.47 & 89.56 & 40.75 & 6.92 & 97.20 & 60.51 & 49.74 & 46.50 \\
AnDPro & 27.97 & \second{44.15} & 55.49 & \second{58.16} & \second{48.19} & 31.92 & 24.80 & 24.50 & \second{24.77} & \second{72.84} & 91.76 & \second{41.99} & 7.31 & \second{99.78} & \best{62.54} & \second{52.00} & 48.01 \\
LAQ++(8) \citep{wang2025lookahead} & 29.64 & 13.22 & 26.79 & 27.58 & 23.49 & 18.63 & 26.94 & 24.04 & 25.21 & 72.50 & 92.25 & 42.83 & 8.43 & 96.25 & 65.00 & 53.37 & 40.39 \\
\textbf{\lorekvE\ (ours)} & \second{29.77} & 43.61 & 55.22 & \best{58.32} & \best{48.68} & \best{32.73} & \second{27.40} & \best{24.96} & \best{25.02} & 72.65 & 91.50 & 40.50 & \second{7.46} & 99.52 & 60.38 & 51.25 & \second{48.06} \\
\textbf{\lorekvD\ (ours)} & \best{30.40} & \best{44.21} & \best{55.87} & 57.31 & 47.46 & \second{32.59} & \best{27.43} & \second{24.91} & 24.68 & \best{73.79} & 92.13 & 40.88 & \best{7.72} & \best{100.00} & 61.40 & \best{53.42} & \best{48.39} \\

\midrule
\multicolumn{18}{c}{\textit{Budget} $=1024$} \\
\midrule
SnapKV & 28.32 & \second{45.37} & \best{56.81} & 58.16 & 48.28 & 31.89 & 26.91 & 24.58 & 25.82 & 70.50 & 92.22 & \second{43.08} & 6.96 & \best{100.00} & \second{62.46} & 51.50 & 48.30 \\
PyramidKV & 29.32 & 44.57 & 56.14 & \second{58.39} & 48.36 & 31.93 & 28.20 & 23.78 & 26.00 & 70.50 & \second{92.41} & 42.46 & 6.88 & \best{100.00} & 61.33 & 49.12 & 48.09 \\
Ada-KV & 28.52 & \best{46.01} & 55.72 & 57.85 & 48.10 & 32.02 & 26.98 & 24.72 & 25.96 & 72.00 & 91.97 & 42.55 & 6.96 & \best{100.00} & 61.98 & 51.52 & 48.30 \\
CriticalKV & 27.27 & 44.74 & 54.24 & 56.17 & 46.95 & 31.65 & 26.35 & 23.91 & 24.93 & 70.22 & 89.08 & 41.00 & 6.70 & 96.86 & 59.85 & 49.75 & 46.85 \\
AnDPro & \best{31.36} & 43.42 & 55.10 & 55.05 & 45.93 & 29.79 & 28.46 & 24.70 & \best{26.56} & 72.50 & 91.42 & \best{43.57} & 7.39 & 99.50 & \best{64.86} & \best{59.01} & 48.66 \\
\textbf{\lorekvE\ (ours)} & 29.99 & 44.83 & 56.19 & \best{58.60} & \best{49.01} & \best{33.02} & \second{28.90} & \best{25.49} & \second{26.38} & \second{73.00} & 92.35 & 42.24 & \second{7.50} & 99.50 & 61.31 & 51.60 & \second{48.74} \\
\textbf{\lorekvD\ (ours)} & \second{30.27} & 45.10 & \second{56.48} & 58.15 & \second{48.46} & \second{32.96} & \best{28.92} & \second{25.47} & 26.23 & \best{73.51} & \best{92.63} & 42.41 & \best{7.62} & \second{99.86} & 61.77 & \second{52.57} & \best{48.90} \\
\bottomrule
\end{tabular}%
}
\end{table}


\begin{table}[p]
\caption{RULER per-task results with Mistral-7B-Instruct-v0.2. CWE/FWE denote common-/frequent-word extraction; MK/MV/MQ/S denote multi-key/multi-value/multi-query/single-needle variants; VT denotes variable tracking.}
\label{tab:ruler_mistral_full}
\centering
\resizebox{\textwidth}{!}{%
\begin{tabular}{llcccccccccccccc}
\toprule
$B$ & Method & CWE & FWE & MK1 & MK2 & MK3 & MQ & MV & S1 & S2 & S3 & QA-1 & QA-2 & VT & Avg. \\
\midrule
\multicolumn{16}{c}{\textit{Context length} $=4$K} \\
\midrule
128 & CriticalKV & 9.40 & 40.20 & 89.20 & 70.40 & \second{0.00} & 0.80 & 17.40 & 89.20 & \second{96.00} & 0.00 & 82.20 & 59.40 & 15.60 & 43.83 \\
128 & AnDPro & 17.20 & 44.20 & \second{98.60} & 87.60 & \best{0.40} & \best{37.40} & \best{46.40} & \second{99.60} & \best{99.80} & 0.00 & 82.60 & 61.20 & 38.00 & 54.85 \\
128 & \textbf{\lorekvE} & \best{29.62} & \best{61.27} & \best{99.80} & \second{99.60} & \second{0.00} & \second{19.40} & \second{36.50} & \best{100.00} & \best{99.80} & \best{4.00} & \second{84.00} & \second{61.60} & \best{63.56} & \best{58.40} \\
128 & \textbf{\lorekvD} & \second{29.18} & \second{60.80} & \best{99.80} & \best{99.80} & \second{0.00} & 19.00 & 36.30 & \best{100.00} & \best{99.80} & \second{3.80} & \best{84.20} & \best{62.00} & \second{63.40} & \second{58.31} \\
\addlinespace[2pt]
256 & CriticalKV & 24.20 & 70.80 & 93.20 & 85.00 & \second{17.20} & 70.00 & 60.20 & 93.20 & 97.40 & 9.00 & 83.20 & 60.80 & 66.40 & 63.89 \\
256 & AnDPro & 29.20 & 73.40 & \second{99.40} & 96.20 & \best{17.40} & \best{93.80} & \best{79.00} & \best{100.00} & \best{99.80} & 9.00 & \second{83.40} & \second{62.00} & 81.00 & \second{71.05} \\
256 & \textbf{\lorekvE} & \second{35.20} & \second{82.40} & 98.60 & \second{98.60} & 15.80 & \second{81.60} & \second{71.80} & \second{98.60} & \second{98.40} & \second{10.00} & 82.80 & 60.80 & \best{95.00} & \best{71.51} \\
256 & \textbf{\lorekvD} & \best{41.52} & \best{84.07} & \best{99.60} & \best{100.00} & 4.80 & 49.15 & 59.25 & \best{100.00} & \best{99.80} & \best{32.80} & \best{84.20} & \best{62.40} & \second{89.44} & 69.77 \\
\addlinespace[2pt]
512 & CriticalKV & 43.60 & 84.00 & \second{95.80} & 92.80 & \second{69.20} & 83.60 & 78.60 & 96.00 & 98.40 & 54.80 & 84.20 & 61.60 & 87.80 & 79.26 \\
512 & AnDPro & 46.40 & 85.40 & \best{99.40} & 99.40 & \best{69.40} & \best{97.40} & \best{89.60} & \best{100.00} & \best{99.80} & 54.80 & \second{84.40} & \best{62.20} & 96.20 & \second{83.42} \\
512 & \textbf{\lorekvE} & \second{50.80} & \best{91.40} & \best{99.40} & \second{99.60} & 68.80 & \second{91.00} & \second{86.00} & \second{99.60} & \second{99.40} & \second{55.80} & \second{84.40} & \second{62.00} & \best{99.80} & \best{83.69} \\
512 & \textbf{\lorekvD} & \best{60.22} & \second{89.53} & \best{99.40} & \best{100.00} & 35.40 & 81.45 & 78.65 & \best{100.00} & \best{99.80} & \best{72.40} & \best{85.40} & 61.80 & \second{97.72} & 81.67 \\
\addlinespace[2pt]
1024 & CriticalKV & 76.40 & 89.00 & 97.60 & 96.60 & \best{89.20} & 90.20 & 84.60 & 98.00 & 99.00 & \second{93.80} & \best{85.00} & 61.20 & 93.60 & 88.78 \\
1024 & AnDPro & 78.00 & 89.80 & \best{99.40} & \best{100.00} & \best{89.20} & \best{97.40} & \best{90.40} & \best{100.00} & \best{99.80} & \second{93.80} & \best{85.00} & \best{61.60} & 98.20 & \second{90.97} \\
1024 & \textbf{\lorekvE} & \second{80.60} & \best{93.20} & \best{99.40} & \second{99.80} & \second{89.00} & 93.80 & \second{88.40} & \second{99.80} & \second{99.60} & \best{94.40} & \best{85.00} & \second{61.40} & \best{100.00} & \best{91.11} \\
1024 & \textbf{\lorekvD} & \best{87.72} & \second{91.53} & \second{99.20} & \best{100.00} & 73.20 & \second{94.00} & 86.70 & \best{100.00} & \best{99.80} & 91.60 & \second{84.80} & 61.00 & \second{98.88} & 89.88 \\
\midrule
\multicolumn{16}{c}{\textit{Context length} $=8$K} \\
\midrule
128 & CriticalKV & 3.20 & 26.40 & 75.20 & 59.60 & \second{0.00} & 0.00 & 6.00 & \second{92.60} & 88.40 & 0.00 & 75.40 & 58.80 & 18.20 & 38.75 \\
128 & AnDPro & 7.40 & 29.60 & 94.60 & 82.80 & \best{0.40} & \best{19.60} & \second{35.40} & \best{100.00} & \second{97.20} & 0.00 & 74.40 & 59.40 & \second{50.00} & 50.06 \\
128 & \textbf{\lorekvE} & \best{19.62} & \best{57.73} & \second{96.40} & \second{98.40} & \second{0.00} & 12.90 & 35.20 & \best{100.00} & \best{100.00} & \best{1.20} & \best{78.20} & \second{61.40} & \best{63.92} & \best{55.77} \\
128 & \textbf{\lorekvD} & \second{19.16} & \second{56.47} & \best{96.80} & \best{98.60} & \second{0.00} & \second{13.10} & \best{35.95} & \best{100.00} & \best{100.00} & \second{0.80} & \second{77.80} & \best{62.00} & \best{63.92} & \second{55.74} \\
\addlinespace[2pt]
256 & CriticalKV & 11.40 & 46.40 & 83.80 & 77.20 & \second{4.00} & 68.80 & 52.20 & 95.20 & 94.00 & \second{2.60} & \best{77.20} & 59.80 & 65.00 & 56.74 \\
256 & AnDPro & 14.20 & 48.40 & \second{96.40} & \second{92.20} & \best{4.20} & \best{81.60} & \best{71.20} & \best{100.00} & \second{99.80} & \second{2.60} & 76.60 & \second{60.20} & 85.60 & 64.08 \\
256 & \textbf{\lorekvE} & \second{19.80} & \second{63.80} & 95.60 & \best{98.20} & 2.20 & \second{75.80} & \second{69.40} & \second{98.20} & 98.20 & 1.60 & \second{77.00} & 59.60 & \best{92.40} & \best{65.52} \\
256 & \textbf{\lorekvD} & \best{29.12} & \best{69.60} & \best{96.80} & \best{98.20} & 1.40 & 34.15 & 55.60 & \best{100.00} & \best{100.00} & \best{25.00} & \second{77.00} & \best{61.40} & \second{91.00} & \second{64.56} \\
\addlinespace[2pt]
512 & CriticalKV & 19.80 & 67.00 & 88.80 & 88.60 & \second{26.60} & 86.60 & 72.40 & 97.20 & 96.60 & \second{19.60} & \best{78.20} & 57.80 & 81.60 & 67.75 \\
512 & AnDPro & 21.40 & 68.20 & \second{96.20} & 97.40 & \best{26.80} & \best{94.00} & \best{83.60} & \best{100.00} & \best{100.00} & \second{19.60} & \second{77.80} & \second{58.00} & 93.80 & 72.06 \\
512 & \textbf{\lorekvE} & \second{24.80} & \best{77.20} & 95.80 & \best{99.00} & 25.60 & \second{90.80} & \second{82.60} & \second{99.00} & \second{99.00} & 19.00 & \best{78.20} & 57.80 & \best{97.80} & \best{72.82} \\
512 & \textbf{\lorekvD} & \best{37.22} & \second{73.07} & \best{96.40} & \second{98.60} & 9.40 & 65.15 & 72.25 & \best{100.00} & \best{100.00} & \best{54.80} & 76.60 & \best{60.40} & \second{94.48} & \second{72.18} \\
\addlinespace[2pt]
1024 & CriticalKV & 64.60 & 80.40 & 94.00 & 94.40 & \best{68.00} & 91.80 & 81.60 & 98.60 & 98.20 & \best{76.20} & \best{81.80} & \second{59.40} & 90.60 & 83.05 \\
1024 & AnDPro & 65.40 & 81.20 & \best{97.80} & 99.00 & \best{68.00} & \best{95.80} & \best{87.40} & \best{100.00} & \best{100.00} & \best{76.20} & \second{81.60} & \best{59.60} & 97.00 & \second{85.31} \\
1024 & \textbf{\lorekvE} & \second{67.60} & \best{86.40} & \second{97.60} & \second{99.40} & \second{67.40} & 94.00 & \second{87.00} & 99.40 & 99.40 & \second{75.80} & \best{81.80} & \second{59.40} & \best{99.40} & \best{85.74} \\
1024 & \textbf{\lorekvD} & \best{74.72} & \second{84.73} & 97.40 & \best{99.60} & 51.60 & \second{94.20} & 85.30 & \second{99.60} & \second{99.60} & 73.00 & \second{81.60} & 59.00 & \second{98.28} & 84.51 \\
\midrule
\multicolumn{16}{c}{\textit{Context length} $=16$K} \\
\midrule
128 & CriticalKV & 0.00 & 44.20 & 62.00 & 50.20 & 0.00 & 0.00 & 0.00 & \second{92.60} & 83.80 & 0.00 & 73.00 & 51.40 & 11.80 & 36.08 \\
128 & AnDPro & 0.00 & 46.40 & 82.00 & 73.80 & 0.00 & 0.00 & \best{24.00} & \best{100.00} & \second{92.80} & 0.00 & 72.20 & 52.40 & 44.00 & 45.20 \\
128 & \textbf{\lorekvE} & \best{12.00} & \best{71.60} & \second{83.60} & \second{88.40} & 0.00 & 0.00 & 21.20 & \best{100.00} & \best{94.80} & \best{1.80} & \best{75.40} & \second{54.00} & \best{60.80} & \best{51.05} \\
128 & \textbf{\lorekvD} & \second{11.55} & \second{70.73} & \best{83.80} & \best{88.60} & 0.00 & 0.00 & \second{21.48} & \best{100.00} & \best{94.80} & \second{1.50} & \second{75.30} & \best{54.50} & \second{60.72} & \second{51.00} \\
\addlinespace[2pt]
256 & CriticalKV & 5.60 & 63.80 & 70.80 & 67.80 & \best{1.40} & \best{49.60} & 44.20 & 95.20 & 89.60 & 1.00 & 75.00 & 52.60 & 58.80 & 51.95 \\
256 & AnDPro & 5.60 & 65.40 & \second{83.80} & 83.20 & \best{1.40} & \best{49.60} & \best{60.00} & \best{100.00} & 95.40 & 1.00 & 74.40 & \second{53.20} & 79.80 & 57.91 \\
256 & \textbf{\lorekvE} & \second{11.80} & \second{79.40} & \second{83.80} & \second{91.00} & \second{0.20} & \second{42.60} & \second{57.20} & \second{99.00} & \second{95.60} & \second{1.20} & \second{75.40} & \second{53.20} & \best{88.80} & \best{59.94} \\
256 & \textbf{\lorekvD} & \best{19.62} & \best{83.14} & \best{84.90} & \best{91.70} & 0.00 & 5.55 & 44.03 & \best{100.00} & \best{97.20} & \best{24.30} & \best{76.10} & \best{54.90} & \second{85.32} & \second{58.98} \\
\addlinespace[2pt]
512 & CriticalKV & 14.60 & 80.80 & 78.60 & 81.20 & \best{24.60} & \best{68.40} & 65.40 & 97.20 & 93.00 & \second{18.40} & \best{76.40} & 52.00 & 76.80 & 63.65 \\
512 & AnDPro & 14.60 & 81.80 & \second{86.20} & \second{90.20} & \best{24.60} & \best{68.40} & \best{74.60} & \best{100.00} & \second{96.40} & \second{18.40} & \second{76.00} & \second{52.40} & 89.00 & \second{67.12} \\
512 & \textbf{\lorekvE} & \second{18.40} & \best{90.20} & 86.00 & \best{94.60} & \second{23.60} & \second{64.60} & \second{73.20} & \second{99.20} & \second{96.40} & 18.20 & \best{76.40} & 52.20 & \best{94.60} & \best{68.28} \\
512 & \textbf{\lorekvD} & \best{22.04} & \second{83.93} & \best{89.00} & 88.00 & 2.80 & 46.60 & 60.70 & \best{100.00} & \best{99.40} & \best{34.20} & \second{76.00} & \best{53.00} & \second{89.04} & 64.98 \\
\addlinespace[2pt]
1024 & CriticalKV & 61.00 & 89.40 & 87.20 & 89.60 & \best{66.60} & \best{79.00} & 76.80 & 98.60 & 95.80 & \best{75.40} & \best{80.60} & \second{55.80} & 87.40 & 80.25 \\
1024 & AnDPro & 61.00 & 89.80 & \best{91.20} & 94.40 & \best{66.60} & \best{79.00} & \best{81.60} & \best{100.00} & \second{97.60} & \best{75.40} & \second{80.40} & \best{56.00} & 93.80 & \second{82.06} \\
1024 & \textbf{\lorekvE} & \second{63.20} & \best{94.60} & \best{91.20} & \second{97.00} & \second{66.20} & 76.80 & \second{80.60} & 99.60 & \second{97.60} & \best{75.40} & \best{80.60} & \best{56.00} & \best{97.00} & \best{82.75} \\
1024 & \textbf{\lorekvD} & \best{70.32} & \second{92.93} & \second{91.00} & \best{97.20} & 50.40 & \second{77.00} & 78.90 & \second{99.80} & \best{97.80} & \second{72.60} & \second{80.40} & 55.60 & \second{95.88} & 81.53 \\
\bottomrule
\end{tabular}%
}
\end{table}\begin{table}[p]
\caption{RULER per-task results with Llama-3.1-8B-Instruct. Abbreviations follow Table~\ref{tab:ruler_mistral_full}.}
\label{tab:ruler_llama_full}
\centering
\resizebox{\textwidth}{!}{%
\begin{tabular}{llcccccccccccccc}
\toprule
$B$ & Method & CWE & FWE & MK1 & MK2 & MK3 & MQ & MV & S1 & S2 & S3 & QA-1 & QA-2 & VT & Avg. \\
\midrule
\multicolumn{16}{c}{\textit{Context length} $=4$K} \\
\midrule
128 & CriticalKV & 9.40 & 22.00 & 91.20 & 78.60 & \second{0.00} & \second{55.80} & 59.40 & \second{91.20} & 96.40 & 0.00 & 86.20 & 58.00 & 11.20 & 50.72 \\
128 & AnDPro & 16.00 & 25.40 & 99.20 & 93.20 & \best{0.20} & \best{86.80} & \best{84.00} & \best{100.00} & \second{99.60} & 0.00 & 86.60 & 59.60 & 30.20 & 60.06 \\
128 & \textbf{\lorekvE} & \best{29.82} & \best{46.67} & \best{100.00} & \second{95.80} & \second{0.00} & 51.55 & 66.95 & \best{100.00} & \best{99.80} & \second{25.40} & \best{88.20} & \second{62.00} & \second{69.36} & \best{64.27} \\
128 & \textbf{\lorekvD} & \second{28.88} & \second{46.07} & \second{99.80} & \best{96.40} & \second{0.00} & 50.55 & \second{67.20} & \best{100.00} & \best{99.80} & \best{25.80} & \second{88.00} & \best{62.60} & \best{69.56} & \second{64.20} \\
\addlinespace[2pt]
256 & CriticalKV & 23.80 & 56.80 & 94.20 & 90.00 & \second{16.80} & 79.40 & 83.60 & 93.80 & 97.20 & 8.80 & \second{85.80} & 59.40 & 62.60 & 65.55 \\
256 & AnDPro & 28.00 & 59.00 & \best{99.40} & \second{99.60} & \best{17.00} & \best{99.60} & \best{99.60} & \best{99.60} & \best{99.20} & 8.80 & \best{86.00} & \best{60.40} & 74.80 & 71.62 \\
256 & \textbf{\lorekvE} & \best{35.00} & \best{69.00} & \second{98.60} & \best{100.00} & 15.60 & \second{86.80} & 92.00 & \second{98.20} & \second{97.80} & \second{10.20} & 85.60 & 59.40 & \second{90.60} & \best{72.22} \\
256 & \textbf{\lorekvD} & \second{34.06} & \second{68.40} & 98.40 & \best{100.00} & 15.60 & 85.80 & \second{92.25} & \second{98.20} & \second{97.80} & \best{10.60} & 85.40 & \second{60.00} & \best{90.80} & \second{72.10} \\
\addlinespace[2pt]
512 & CriticalKV & 43.00 & 74.20 & 96.00 & 93.60 & \second{68.60} & 87.40 & 89.80 & 95.80 & 97.60 & 54.20 & 85.40 & 60.20 & 84.40 & 79.25 \\
512 & AnDPro & 45.40 & 75.40 & \second{98.80} & 99.00 & \best{68.80} & \best{99.00} & \best{99.00} & \best{99.00} & \best{98.60} & 54.20 & \second{85.60} & \second{60.80} & 91.60 & 82.71 \\
512 & \textbf{\lorekvE} & \best{49.80} & \best{81.40} & \best{99.00} & \second{99.40} & 68.40 & \second{92.60} & 95.40 & \second{98.80} & \second{98.40} & \second{55.40} & \best{85.80} & \second{60.80} & \second{95.40} & \best{83.12} \\
512 & \textbf{\lorekvD} & \second{48.86} & \second{80.80} & \second{98.80} & \best{100.00} & 68.40 & 91.60 & \second{95.65} & \second{98.80} & \second{98.40} & \best{55.80} & \second{85.60} & \best{61.40} & \best{95.60} & \second{83.05} \\
\addlinespace[2pt]
1024 & CriticalKV & 75.40 & 82.40 & 96.80 & 96.00 & 88.00 & 92.60 & 93.80 & 97.00 & 98.00 & 92.60 & 85.00 & 60.00 & 91.00 & 88.35 \\
1024 & AnDPro & 86.00 & 89.60 & \best{99.80} & \best{100.00} & \best{98.40} & \best{99.80} & \best{99.80} & \best{100.00} & \best{100.00} & 95.60 & \best{87.80} & \second{62.60} & 94.60 & 93.38 \\
1024 & \textbf{\lorekvE} & \best{88.80} & \best{93.40} & \best{99.80} & \second{99.80} & \second{98.20} & \second{96.20} & 97.80 & \second{99.80} & \second{99.80} & \second{96.40} & \best{87.80} & 62.40 & \second{96.60} & \best{93.60} \\
1024 & \textbf{\lorekvD} & \second{87.86} & \second{92.80} & \second{99.60} & \best{100.00} & \second{98.20} & 95.20 & \second{98.05} & \second{99.80} & \second{99.80} & \best{96.80} & \second{87.60} & \best{63.00} & \best{96.80} & \second{93.50} \\
\midrule
\multicolumn{16}{c}{\textit{Context length} $=8$K} \\
\midrule
128 & CriticalKV & 8.20 & 13.80 & 82.80 & 73.80 & \second{0.00} & \second{68.20} & 58.60 & 93.60 & 92.00 & 0.00 & 81.80 & 59.40 & 2.80 & 48.85 \\
128 & AnDPro & 11.80 & 16.60 & 99.20 & \second{93.60} & \second{0.00} & \best{84.80} & 83.40 & \second{99.80} & \second{99.40} & 0.00 & 81.00 & 60.00 & 29.80 & 58.42 \\
128 & \textbf{\lorekvE} & \best{33.20} & \best{39.00} & \best{100.00} & \best{100.00} & \best{8.20} & 36.60 & \second{87.20} & \best{100.00} & \best{100.00} & \second{9.40} & \best{94.40} & \second{70.80} & \second{45.80} & \best{63.43} \\
128 & \textbf{\lorekvD} & \second{32.26} & \second{38.40} & \second{99.80} & \best{100.00} & \best{8.20} & 35.60 & \best{87.45} & \best{100.00} & \best{100.00} & \best{9.80} & \second{94.20} & \best{71.40} & \best{46.00} & \second{63.32} \\
\addlinespace[2pt]
256 & CriticalKV & 15.00 & 36.80 & 88.60 & 86.80 & \second{3.60} & 88.60 & 83.40 & 95.40 & 94.60 & \best{2.40} & \best{81.80} & 59.80 & 52.80 & 60.74 \\
256 & AnDPro & 17.20 & 38.60 & \best{99.40} & \second{99.60} & \best{3.80} & \best{99.60} & \best{99.60} & \best{99.40} & \best{99.60} & \best{2.40} & 81.20 & \second{60.20} & 70.20 & 66.98 \\
256 & \textbf{\lorekvE} & \best{23.80} & \best{55.80} & \second{98.60} & \best{100.00} & 2.00 & \second{93.60} & 97.80 & \second{97.60} & \second{98.00} & 1.60 & \best{81.80} & 59.80 & \second{78.00} & \best{68.34} \\
256 & \textbf{\lorekvD} & \second{22.86} & \second{55.20} & 98.40 & \best{100.00} & 2.00 & 92.60 & \second{98.05} & \second{97.60} & \second{98.00} & \second{2.00} & \second{81.60} & \best{60.40} & \best{78.20} & \second{68.22} \\
\addlinespace[2pt]
512 & CriticalKV & 22.00 & 60.20 & 91.40 & 91.80 & \second{26.00} & 92.80 & 89.60 & 96.80 & 96.40 & \best{19.40} & \second{80.60} & 57.40 & 72.80 & 69.02 \\
512 & AnDPro & 23.40 & 61.20 & \best{97.60} & \second{99.00} & \best{26.20} & \best{99.00} & \best{99.00} & \best{99.00} & \best{99.00} & \best{19.40} & 80.20 & \second{57.60} & 83.00 & 72.58 \\
512 & \textbf{\lorekvE} & \best{27.00} & \best{70.60} & \second{97.20} & \best{100.00} & 25.20 & \second{95.80} & 98.20 & \second{98.20} & \second{98.20} & \second{19.00} & \best{80.80} & \second{57.60} & \second{87.20} & \best{73.46} \\
512 & \textbf{\lorekvD} & \second{26.06} & \second{70.00} & 97.00 & \best{100.00} & 25.20 & 94.80 & \second{98.45} & \second{98.20} & \second{98.20} & \best{19.40} & \second{80.60} & \best{58.20} & \best{87.40} & \second{73.35} \\
\addlinespace[2pt]
1024 & CriticalKV & \best{65.20} & 75.80 & 94.80 & 94.80 & 66.80 & 95.40 & 93.80 & 97.60 & 97.20 & 75.20 & 82.60 & 58.80 & 84.40 & 83.26 \\
1024 & AnDPro & 45.00 & 81.60 & \best{100.00} & \second{99.80} & \best{76.40} & \best{99.60} & \best{99.80} & \best{100.00} & \best{100.00} & \best{77.20} & 83.40 & \second{61.80} & 92.60 & 85.94 \\
1024 & \textbf{\lorekvE} & \second{47.20} & \best{87.20} & \second{99.80} & \best{100.00} & \second{75.80} & \second{97.80} & 99.20 & \second{99.40} & \second{99.40} & \second{76.80} & \best{83.80} & 61.60 & \second{95.20} & \best{86.40} \\
1024 & \textbf{\lorekvD} & 46.26 & \second{86.60} & 99.60 & \best{100.00} & \second{75.80} & 96.80 & \second{99.45} & \second{99.40} & \second{99.40} & \best{77.20} & \second{83.60} & \best{62.20} & \best{95.40} & \second{86.29} \\
\midrule
\multicolumn{16}{c}{\textit{Context length} $=16$K} \\
\midrule
128 & CriticalKV & 2.80 & 30.00 & 68.60 & 63.20 & \second{0.00} & 61.20 & 49.00 & \second{93.80} & 86.80 & 0.00 & 78.80 & 51.60 & 0.00 & 45.06 \\
128 & AnDPro & 2.80 & 32.00 & 85.60 & 83.20 & \second{0.00} & 61.20 & 69.60 & \best{100.00} & \second{94.40} & 0.00 & 78.20 & 52.40 & 27.00 & 52.80 \\
128 & \textbf{\lorekvE} & \best{4.20} & \best{34.60} & \best{90.80} & \second{88.40} & \best{1.20} & \best{65.40} & \second{74.00} & \best{100.00} & \best{100.00} & \second{1.20} & \best{83.20} & \second{56.00} & \second{29.60} & \best{56.05} \\
128 & \textbf{\lorekvD} & \second{3.26} & \second{34.00} & \second{90.60} & \best{89.00} & \best{1.20} & \second{64.40} & \best{74.25} & \best{100.00} & \best{100.00} & \best{1.60} & \second{83.00} & \best{56.60} & \best{29.80} & \second{55.98} \\
\addlinespace[2pt]
256 & CriticalKV & 7.60 & 53.00 & 75.00 & 77.00 & \best{1.20} & \best{95.00} & 80.40 & 95.60 & 91.20 & 1.00 & 79.00 & 52.20 & 49.00 & 58.25 \\
256 & AnDPro & 7.60 & 54.20 & \second{86.00} & 90.00 & \best{1.20} & \best{95.00} & \best{93.60} & \best{99.60} & \second{96.20} & 1.00 & 78.40 & 52.80 & 66.60 & 63.25 \\
256 & \textbf{\lorekvE} & \best{14.80} & \best{70.20} & \best{86.20} & \second{99.00} & \second{0.00} & \second{88.00} & 91.00 & \second{98.60} & \best{96.60} & \second{1.40} & \best{79.80} & \second{53.00} & \second{76.80} & \best{65.80} \\
256 & \textbf{\lorekvD} & \second{13.86} & \second{69.60} & \second{86.00} & \best{99.60} & \second{0.00} & 87.00 & \second{91.25} & \second{98.60} & \best{96.60} & \best{1.80} & \second{79.60} & \best{53.60} & \best{77.00} & \second{65.73} \\
\addlinespace[2pt]
512 & CriticalKV & 16.00 & 73.20 & 80.80 & 86.60 & \best{24.20} & \best{98.20} & 88.80 & 96.80 & 93.60 & \second{18.20} & \second{78.60} & 51.40 & 69.60 & 67.38 \\
512 & AnDPro & 16.00 & 74.00 & \best{87.20} & 94.20 & \best{24.20} & \best{98.20} & \best{96.40} & \best{99.00} & \second{96.40} & \second{18.20} & 78.20 & \second{51.80} & 79.80 & 70.28 \\
512 & \textbf{\lorekvE} & \best{20.00} & \best{82.60} & \best{87.20} & \second{98.80} & \second{23.20} & \second{94.40} & 95.20 & \second{98.40} & \best{96.60} & \second{18.20} & \best{78.80} & \second{51.80} & \second{85.60} & \best{71.60} \\
512 & \textbf{\lorekvD} & \second{19.06} & \second{82.00} & \second{87.00} & \best{99.40} & \second{23.20} & 93.40 & \second{95.45} & \second{98.40} & \best{96.60} & \best{18.60} & \second{78.60} & \best{52.40} & \best{85.80} & \second{71.53} \\
\addlinespace[2pt]
1024 & CriticalKV & \best{61.20} & \best{84.20} & 87.60 & 92.00 & \best{65.40} & 96.20 & 90.00 & 97.60 & \second{95.20} & \best{74.40} & \best{81.40} & 55.00 & 82.20 & \best{81.72} \\
1024 & AnDPro & 24.80 & 75.80 & \best{99.80} & \second{99.40} & \second{57.80} & \best{99.40} & \best{98.80} & \best{100.00} & \best{100.00} & 51.00 & 80.80 & \second{55.80} & 88.40 & 79.37 \\
1024 & \textbf{\lorekvE} & \second{27.20} & \second{81.00} & \best{99.80} & \best{100.00} & 57.40 & \second{97.00} & 98.00 & \second{99.60} & \best{100.00} & 51.00 & \second{81.20} & \second{55.80} & \second{91.80} & \second{79.98} \\
1024 & \textbf{\lorekvD} & 26.26 & 80.40 & \second{99.60} & \best{100.00} & 57.40 & 96.00 & \second{98.25} & \second{99.60} & \best{100.00} & \second{51.40} & 81.00 & \best{56.40} & \best{92.00} & 79.87 \\
\bottomrule
\end{tabular}%
}
\end{table}

\subsection{LongBench with Llama-3.1-8B-Instruct}
\label{app:llama_results}

Table~\ref{tab:longbench_llama} provides the per-dataset Llama-3.1-8B-Instruct results. CriticalKV, AnDPro, and \lorekvE\ are directly comparable because they use the same implementation, whereas SnapKV, PyramidKV, and Ada-KV are published results from \citet{geng2026accurate}. \lorekvE\ has the highest 16-task average at all four budgets, although its margin over AnDPro is small; the \lorekvD\ rows are \tbd\ at every budget and use the same implementation and particles, differing only in $(\gamma,\lambda_D)$.



At $B{=}128$, \lorekvE\ shows the largest gains over AnDPro on multi-hop and long-document QA, with smaller improvements on summarization, while performance drops mainly on code completion and a few retrieval/few-shot tasks. This trend is consistent with the per-dataset results in Table~\ref{tab:longbench_llama}; the corresponding Mistral code results in Table~\ref{tab:longbench_mistral} are closer to neutral. A separate $M{=}1$, $T{=}8$ Llama run was completed only on NarrativeQA, Qasper, and MultiFieldQA-en and is therefore excluded from Table~\ref{tab:longbench_llama}. Lookahead-count comparisons for Llama are instead based on the matched RULER ablation in Table~\ref{tab:ruler_ablation}.

\subsection{RULER: Protocol and Per-Task Results} \label{app:ruler_results}

All RULER results are obtained under one matched implementation using chunk size $4$, $w{=}32$, and $B\in\{128,256,512,1024\}$; \lorekvE\ follows the core configuration in Table~\ref{tab:hyperparams}, and \lorekvD\ differs from it only in $(\gamma,\lambda_D)=(1,0.3)$. Each task is scored on its native $0$--$100$ metric, averaged across the 13 tasks at each context length and then across $4$K, $8$K, and $16$K; Tables~\ref{tab:ruler_mistral_full} and~\ref{tab:ruler_llama_full} report the complete per-task results. At tight budgets, \lorekvE\ gains are strongest on aggregation-style tasks such as VT, FWE, and CWE, while single-needle tasks become less discriminative as the budget increases. The main exception is the multi-query needle task, where \lorekvE\ can underperform AnDPro at shorter contexts, and NIAH-MK3 remains difficult for all methods at $B{=}128$. At the opposite extreme, CriticalKV leads the Llama $B{=}1024$, $16$K setting ($81.72$ versus $79.98$ for \lorekvE\ and $79.37$ for AnDPro), which we report as an isolated exception to the overall trend.


\subsection{RULER Ablations} \label{app:ruler_ablation}
Table~\ref{tab:ruler_ablation} tests the main ablation axes on RULER, whose prompt structure differs from LongBench. The cleanest comparison is at $4$K and $B{=}128$, where $M{=}1$ and the \lorekvE\ core use matched uniform head budgets and top-$K$ selection: increasing $M$ from $1$ to $4$ improves the average from $56.90$ to $58.43$. By comparison, moving $\gamma$ from $1$ to $1.3$ changes the corresponding $M{=}4$ score from $58.39$ to $58.43$---a $23\%$ perturbation along a single ray of the $(\beta,\gamma)$ plane that never reaches $\gamma{=}0$, so it bounds the effect of that ray and not of the weighting apparatus, and the adaptive-budget and diversity switches each change the score by at most $0.11$ in the available matched rows. At $B{=}256$, the available adaptive-budget comparison similarly improves from $68.50$ with $M{=}1$ to $70.04$ with $M{=}4$. Most variants were evaluated only at $4$K. In particular, the available $M{=}1$ result at $16$K uses adaptive rather than uniform head budgets, so the $52.18$ versus $51.05$ comparison is not a single-factor lookahead-count ablation.

\textbf{Complete Configuration Grid} Table~\ref{tab:ablation_summary} reports all executed LongBench configurations with Mistral-7B-Instruct-v0.2. It complements Table~\ref{tab:ablation_summary} by exposing the exact hyperparameter vector of every run; rows that change multiple switches are included for completeness but are not interpreted as single-factor ablations.

\begin{table}[t]
\caption{LongBench ablations (16-dataset average) with
Mistral-7B-Instruct-v0.2. \tbd\ denotes configurations not yet run.}
\label{tab:ablation_summary}
\centering
\setlength{\tabcolsep}{3.5pt}
\renewcommand{\arraystretch}{1.12}
\resizebox{\columnwidth}{!}{%
\begin{tabular}{@{}llcccc@{}}
\toprule
\textbf{Axis} & \textbf{Variant}
& $B{=}128$ & $B{=}256$ & $B{=}512$ & $B{=}1024$ \\
\midrule

\multirow{2}{*}{Core}
& \lorekvE\ ($M{=}4$, $(1,1.3,0)$)
& 40.28 & 41.38 & 41.86 & 42.20 \\
& \lorekvD\ ($M{=}4$, $(1,1,0.3)$)
& 40.39 & 41.44 & 41.86 & 42.17 \\
\midrule

\multirow{4}{*}{Lookahead budget}
& $M{=}0$ (prompt window only; $\lambda_L{=}0$)
& 35.87 & 37.50 & 38.70 & 39.50 \\
& $M{=}1$ (single surrogate response)
& 39.81 & 40.85 & 41.56 & 42.03 \\
& $M{=}2$
& 40.17 & 40.18 & 41.10 & 41.22 \\
& $M{=}8$
& 40.69 & 40.97 & 41.52 & 41.65 \\
\midrule

\multirow{6}{*}{Reliability weighting}
& Uniform weights, $(0,0,0)$
& \tbd & \tbd & \tbd & \tbd \\
& Log-probability only, $(1,1,0)$
& 40.29 & 41.39 & 41.30 & 41.40 \\
& Pure typicality, $(1,0,0)$
& \tbd & \tbd & \tbd & \tbd \\
& Entropy $+$ disagreement, $(1,1.3,0.3)$
& 40.33 & 41.40 & 41.84 & 42.17 \\
& Disagreement without entropy, $(1,0,0.3)$
& \tbd & \tbd & \tbd & \tbd \\
& Stronger entropy, $(1,2,0)$
& \tbd & \tbd & \tbd & \tbd \\
\midrule

\multirow{3}{*}{Perturbation objective}
& Accumulated attention mass
& \tbd & \tbd & \tbd & \tbd \\
& Unprojected norm, $\|\Delta\mathbf{o}\|_2^2$
& \tbd & \tbd & \tbd & \tbd \\
& AnDPro scoring with the same continuations
& \tbd & \tbd & \tbd & \tbd \\
\midrule

\multirow{2}{*}{Plug-ins}
& $+$ Adaptive head budgets
& 40.28 & 41.35 & 41.72 & 41.72 \\
& $+$ Adaptive budgets and diversity
& 40.30 & 40.33 & 41.50 & 41.80 \\
\bottomrule
\end{tabular}%
}

\vspace{2pt}
\parbox{\columnwidth}{\footnotesize
\textit{Configuration details.}
Unless otherwise stated, $T{=}8$, $\lambda_L{=}1$, uniform head
budgets, and top-$K$ selection are used.
Adaptive head budgets use $\tau{=}1$, while diversity adds
$\lambda_k{=}\lambda_v{=}0.1$.
Reliability weighting is vacuous at $M{=}1$, since the softmax over a
single trajectory always gives unit weight.
For $M{=}2$, $\hat D_m$ is identically zero by symmetry, so trajectory
disagreement cannot affect the weights, although the typicality and
entropy coefficients can.
$M{=}0$ denotes the prompt-window-only configuration.}
\end{table}


\begin{table}[t]
\caption{RULER ablations with Mistral-7B-Instruct-v0.2
(13-task averages).}
\label{tab:ruler_ablation}
\centering
\setlength{\tabcolsep}{4.5pt}
\resizebox{\textwidth}{!}{%
\begin{tabular}{@{}lccccccccccc@{}}
\toprule
& & & & & &
\multicolumn{3}{c}{$B{=}128$}
& \multicolumn{3}{c}{$B{=}256$} \\
\cmidrule(lr){7-9}\cmidrule(lr){10-12}

\textbf{Configuration}
& $M$ & $\gamma$ & $\lambda_D$
& \textbf{Head budgets} & \textbf{Selection}
& 4K & 8K & 16K
& 4K & 8K & 16K \\
\midrule

AnDPro (strongest baseline)
& -- & -- & -- & -- & --
& 54.85 & 50.06 & 45.20
& 71.05 & 64.08 & 57.91 \\

\midrule

$M{=}1$, uniform, top-$K$
& 1 & 1 & 0 & uniform & top-$K$
& 56.90 & 54.29 & 52.15
& 68.45 & 62.77 & 57.70 \\

$M{=}1$, adaptive, top-$K$
& 1 & 1 & 0 & adaptive & top-$K$
& 56.92 & 54.31 & 52.18
& 68.50 & 62.81 & 57.75 \\

$M{=}1$, adaptive, diversity
& 1 & 1 & 0 & adaptive & diversity
& 56.95 & 54.34 & 52.20
& 68.55 & 62.85 & 57.78 \\

\midrule

$M{=}4$, $(1,1,0)$, uniform, top-$K$
& 4 & 1 & 0 & uniform & top-$K$
& 58.39 & 55.72 & 50.98
& 70.08 & 64.45 & 58.85 \\

$M{=}4$, $(1,1,0)$, adaptive, top-$K$
& 4 & 1 & 0 & adaptive & top-$K$
& 58.28 & 55.68 & 50.95
& 70.04 & 64.42 & 58.82 \\

$M{=}4$, $(1,1,0)$, adaptive, diversity
& 4 & 1 & 0 & adaptive & diversity
& 58.31 & 55.70 & 50.97
& 70.07 & 64.45 & 58.85 \\

$M{=}4$, adaptive, top-$K$
& 4 & 1.3 & 0 & adaptive & top-$K$
& 58.37 & 55.74 & 51.02
& 71.45 & 65.46 & 59.88 \\

\textbf{$M{=}4$, uniform, top-$K$ (\lorekvE\ core)}
& 4 & 1.3 & 0 & uniform & top-$K$
& \textbf{58.40} & \textbf{55.77} & 51.05
& \textbf{71.51} & \textbf{65.52} & \textbf{59.94} \\

\midrule

$M{=}4$, disagreement, adaptive, top-$K$
& 4 & 1 & 0.3 & adaptive & top-$K$
& 58.28 & 55.70 & 50.96
& 69.72 & 64.50 & 58.92 \\

$M{=}4$, disagreement, adaptive, diversity
& 4 & 1 & 0.3 & adaptive & diversity
& 58.31 & 55.72 & 50.98
& 69.75 & 64.53 & 58.95 \\

\textbf{$M{=}4$, uniform, top-$K$ (\lorekvD\ core)}
& 4 & 1 & 0.3 & uniform & top-$K$
& 58.31 & 55.74 & 51.00
& 69.77 & 64.56 & 58.98 \\

\bottomrule
\end{tabular}%
}
\end{table}



\subsection{Needle-in-a-Haystack Grid}
\label{app:niah_extra}
\label{app:large_models}

Figure~\ref{fig:niah-heatmap} gives the complete depth$\times$length grid for Llama-3.1-8B-Instruct at $B{=}1024$, with \lorekvE\ as the LORE-KV panel; the \lorekvD\ grid is pending. Retrieval is close to saturation at this budget, with differences concentrated at the more difficult depth--length combinations, so this setting has little power to separate the two reliability variants---which is why the \lorekvE-versus-\lorekvD\ comparison of Appendix~\ref{app:reparam_grid} runs at $B{=}128$ instead. RULER provides the controlled low-budget counterpart used in the main analysis.

\begin{figure}[htbp]
  \centering
  \begin{subfigure}[t]{0.49\textwidth}
    \centering
    \includegraphics[width=\textwidth]{figures_combined/niah_heatmap_exact_match_llama_b1024.pdf}
    \subcaption{Exact match}
    \label{fig:niah-em}
  \end{subfigure}
  \hfill
  \begin{subfigure}[t]{0.49\textwidth}
    \centering
    \includegraphics[width=\textwidth]{figures_combined/niah_heatmap_recall_llama_b1024.pdf}
    \subcaption{Word recall}
    \label{fig:niah-recall}
  \end{subfigure}
  \caption{\textbf{Needle-in-a-Haystack retrieval.} Exact-match performance across needle depths and 4K--32K context lengths with Llama-3.1-8B-Instruct, $B=1024$.}
  \label{fig:niah-heatmap}
\end{figure}




\section{Mechanism Analysis}
\label{app:experiments_extra}

\subsection{Measurement Protocol}
\label{app:mechanism_protocol}

All mechanism analyses use a common capture procedure applied to the actual eviction implementation.

\textbf{Captured states.}
For each prompt, we (i) run the standard prefill and retain the post-RoPE prompt-window queries and full
$\mathbf{K}^{l,h},\mathbf{V}^{l,h}$; (ii) sample the $M$ lookahead trajectories and store their query states, the per-step centered score $c^{(m)}_u$, entropy $H^{(m)}_u$ and varentropy $V^{(m)}_u$, the raw and standardized trajectory disagreement $D_m,\hat D_m$, and the weight $\omega_m$, with the three channels of Eq.~\eqref{eq:confidence} stored separately; per-step storage costs $MT$ floats per prompt; and (iii) greedy-decode against the \emph{full} cache and record the post-RoPE queries at each decoding step. Disagreement is recorded even when $\lambda_D{=}0$, so that both reliability signals are available as diagnostics on every capture regardless of which variant produced it; the analysis code imports the same $D_m$ routine the deployed method uses, so the figures and the runtime score identically by construction. Step~(iii) defines $\mathcal{Q}_{\mathrm{future}}$ and is used only for analysis. Queries are recomputed from each layer's own projection and rotary embedding, leaving the forward pass unchanged.

\textbf{Analysis score.}
All token-importance analyses use the LORE-KV scoring rule. For query $\mathbf{q}$ and candidate token $i$, let
$c_{q,i}=a^{l,h}_{q,i}/(1-a^{l,h}_{q,i}+\epsilon)$. We use
\begin{equation}
s^{l,h}_{q,i}
=
\left|c_{q,i}\right|
\left\|
\mathbf{W}^{O}_{l,h}\mathbf{v}^{l,h}_{i}
-
\mathbf{W}^{O}_{l,h}\mathbf{o}^{l,h}_{q}
\right\|_2 .
\label{eq:analysis_score}
\end{equation}
Scores are averaged over the corresponding query set and combined using Eq.~\eqref{eq:final_score}. Prompt-window queries retain their causal mask; lookahead and future queries occur after the full prompt and require none.

\textbf{Compared criteria.}
All criteria use the same captured model state and selector. \textbf{Attention} is the mean SnapKV-style prompt-window attention mass; \textbf{AnDPro window} uses the projected perturbation with $\lambda_L{=}0$; \textbf{LORE-KV} uses Eq.~\eqref{eq:final_score} with the \lorekvE\ weights unless a panel states otherwise; and \textbf{Single-lookahead} combines the window with one unweighted trajectory as a score-level proxy for Lookahead Q-Cache. Protected positions $\mathcal{P}$ receive score $+\infty$, and the remaining tokens are selected by uniform per-head top-$\min(B,n)$. Thus, comparisons isolate the scoring criterion rather than downstream selection.

\textbf{Analysis-only deviations.}
Four deviations from deployment apply to these figures. \textbf{D1:} Eq.~\eqref{eq:analysis_score} uses the magnitude rather than squared magnitude of Eq.~\eqref{eq:projected_score}; this preserves single-query rankings but changes the scale of the layer energy in Figure~\ref{fig:layer_head_analysis}. \textbf{D2:} chunking, adaptive head budgets, and diversity are disabled, and $B$ includes the ${\sim}33$ protected positions. \textbf{D3:} Figure~\ref{fig:layer_head_analysis} sums energy over all positions and displays GQA-expanded attention heads rather than KV heads. \textbf{D4:} Figure~\ref{fig:retained_token_composition} identifies answer tokens by exact token-id match to the gold answer; this may underestimate absolute evidence retention but affects all criteria equally.

\textbf{Derived quantities.}
Coverage measures how closely an estimated query set matches the true future queries:
\begin{equation}
\operatorname{Cov}\big(\mathcal{Q}_{\mathrm{future}};
\{\hat{\mathcal{Q}}_s\},\boldsymbol{\pi}\big)
=
\frac{1}{H}\sum_h
\frac{1}{T_{\mathrm f}}\sum_t
\sum_s \pi_s
\max_{\hat{\mathbf q}\in\hat{\mathcal Q}_s}
\cos\!\big(
\mathbf q^{\mathrm{future}}_{t,h},
\hat{\mathbf q}
\big).
\label{eq:coverage}
\end{equation}
Queries are unit-normalized before computing cosine similarity. The oracle score $S^\star$ is Eq.~\eqref{eq:final_score} evaluated directly on $\mathcal{Q}_{\mathrm{future}}$, and
\begin{equation}
\operatorname{Recall}@B(S,S^\star)
=
\frac{
\left|\mathrm{top}_B(S)\cap\mathrm{top}_B(S^\star)\right|
}{B}
\label{eq:recall_at_b}
\end{equation}
measures agreement with oracle retention. Behavioral fidelity evaluates the retained cache directly:
\begin{equation}
F(B)
=
\mathbb{E}_{\mathbf q\sim\mathcal Q_{\mathrm{future}}}
\cos\!\big(
\mathbf o^{\mathcal R(B)}_q,
\mathbf o_q
\big).
\label{eq:fidelity}
\end{equation}
Finally,
$\operatorname{Mis}(X)=1-\operatorname{Cov}(\mathcal{Q}_{\mathrm{future}};
\{\mathcal Q_{\mathcal W}^{l,h}\})$
denotes prompt-to-future query mismatch.

\textbf{Configuration.}
Unless noted in Table~\ref{tab:figure_recipes}, all mechanism analyses use Mistral-7B-Instruct-v0.2, the core configuration of Table~\ref{tab:hyperparams} with the entropy weighting of \lorekvE, $B{=}128$, LongBench single-document QA, and RULER at $4$K. The two reliability variants are compared explicitly only in Figures~\ref{fig:future_query_alignment} and~\ref{fig:particle_divergence}, which re-weight the captured particles under both rules; every other figure is insensitive to the choice, because it either fixes the weights or compares criteria that do not use them. Scatter plots may thin displayed markers, but all statistics use the complete samples.

\begin{table}[t]
\caption{Mechanism-analysis measurements, lookahead settings, and sample sizes.}
\label{tab:figure_recipes}
\centering
\setlength{\tabcolsep}{4pt}
\resizebox{\textwidth}{!}{%
\begin{tabular}{@{}lllccl@{}}
\toprule
Figure & Grouping variable & Measured quantity & Steps & $M$ & Sample size (LB / RULER) \\
\midrule
\ref{fig:future_query_alignment} & particles $M\in\{1,2,4,8\}$ & coverage, Eq.~\eqref{eq:coverage} & 32 & 8 & 118 / 200 prompts \\
\ref{fig:score_comparison} & candidate rank; criterion pair & score profile; overlap; Spearman $\rho$ & 0 & 4 & 4{,}501 / 3{,}843 candidates \\
\ref{fig:cache_distribution} & cosine value; layer index & fidelity, Eq.~\eqref{eq:fidelity} & 16 & 4 & 2.47M / 3.28M triples \\
\ref{fig:confidence_calibration} & binned confidence $R_m$; kept fraction & $\operatorname{Recall}@B$; $1-\operatorname{Cov}$ & 32 & 8 & 520 / 1{,}600 particles \\
\ref{fig:mismatch_vs_gain} & mismatch $\operatorname{Mis}(X)$ & recall and fidelity gain & 24 & 4 & 265 / 200 examples \\
\ref{fig:retained_token_composition} & token class & retention; enrichment & 0 & 4 & 751K / 769K tokens \\
\ref{fig:layer_head_analysis} & layer; (layer, head) & energy $E_l$; recall gain $G_{l,h}$ & 24 & 4 & $32\times32$ grid \\
\ref{fig:budget_equivalence} & budget $B\in\{64,\ldots,1024\}$ & fidelity $F(B)$; implied saving & 24 & 4 & one score pass, five selections \\
\ref{fig:particle_divergence} & disagreement $D_m$; entropy $\bar H_m$; signal; weighting rule & $\operatorname{Recall}@B$; Spearman $\rho$; coverage, Eq.~\eqref{eq:coverage} & 32 & 8 & 4{,}400 / 52{,}000 particles \\
\bottomrule
\end{tabular}%
}
\end{table}


\subsection{Future-Query Coverage}
\label{app:query_coverage}

Figure~\ref{fig:future_query_alignment} compares the prompt-end window, one unweighted rollout, uniformly weighted multi-trajectory queries, and the two reliability-weighted estimators of Eq.~\eqref{eq:particle_weight} using Eq.~\eqref{eq:coverage}. The $M$-sweep uses nested prefixes of the same eight-particle sample, so all multi-trajectory variants coincide at $M{=}1$. Because $\hat D_m$ is ensemble-relative, it is recomputed within each nested prefix rather than sliced from the full-ensemble value, which keeps the \lorekvD\ curve a valid estimator at every $M$. Response-side queries produce the dominant coverage increase, while additional particles give diminishing returns after roughly $M{=}4$, and the two weighting rules separate only slightly---the ordering of the weights matters far less than the decision to use response-side queries at all. The smaller separation on RULER is consistent with its question-at-prompt-end format.

\begin{figure}[htbp]
  \centering
  \includegraphics[width=\textwidth]{figures_combined/fig1_future_query_alignment_update.pdf}
  \caption{\textbf{Future-query coverage.} Prompt-window, single-trajectory, and multi-trajectory estimates as the number of lookahead particles increases. Bands show $\pm1$ s.e.m.}
  \label{fig:future_query_alignment}
\end{figure}


\subsection{Retention-Order Shift}
\label{app:score_shift}

Figure~\ref{fig:score_comparison} measures whether lookahead changes which tokens cross the retention boundary rather than merely rescaling prompt-side scores. At $B{=}128$, attention and \lorekvE\ have top-$B$ overlap $0.82$ on LongBench and $0.86$ on RULER, despite high rank correlations ($\rho=0.94$ and $0.87$). Thus, most of the global ordering is shared while the tokens near the budget boundary differ. Single-lookahead and \lorekvE\ are substantially closer ($0.97$ overlap and $\rho=1.00$ on LongBench; $0.92$ and $0.95$ on RULER), providing the score-level compute-matched comparison with a single response-side surrogate. The corresponding \lorekvE-versus-\lorekvD\ overlap on the same particles gates whether an end-to-end run can resolve the two variants at all.

\begin{figure}[t]
  \centering
  \includegraphics[width=\textwidth]{figures_combined/fig2_score_comparison.pdf}
  \caption{\textbf{Retention-order shift.} Candidate-score profiles, top-$B$ overlap, and Spearman correlation across eviction criteria. Mistral-7B-Instruct-v0.2, $B=128$, $M=4$.}
  \label{fig:score_comparison}
\end{figure}


\subsection{Particle-Confidence Calibration}
\label{app:calibration_extra}
\label{app:confidence_details}

Figure~\ref{fig:confidence_calibration} evaluates whether the combined confidence $R_m$, at the \lorekvE\ coefficients $(1,1.3,0)$, predicts oracle retention utility. Utility increases weakly with confidence, while hard removal of low-confidence particles provides little benefit: the residual coverage error is essentially flat in the kept fraction, and retaining only the most confident particles does not beat keeping all of them. This supports soft weighting rather than hard particle selection and is consistent with the small difference between $\gamma{=}1$ and the core $\gamma{=}1.3$ in the ablations---though, as Section~\ref{subsec:confidence} notes, that pair brackets a narrow range of one coefficient rather than the reachable set. Which individual signal carries this weak correlation---average log-probability, $-\bar H_m$, $-\hat D_m$, or either combined score---is measured in the third row of Figure~\ref{fig:particle_divergence}, which ranks all of them against the same utility on one set of particles; the same figure supplies the entropy-versus-utility relationship itself, so it is not repeated here.

\begin{figure}[t]
  \centering
  \includegraphics[width=\textwidth]{figures_combined/fig5_confidence_calibration.pdf}
  \caption{\textbf{Particle-confidence calibration.} Oracle retention recall versus reliability score (top) and residual query-coverage error after confidence filtering (bottom).}
  \label{fig:confidence_calibration}
\end{figure}


\subsection{Trajectory Disagreement as the Reliability Signal}
\label{app:divergence}

Figure~\ref{fig:particle_divergence} prices each channel of Eq.~\eqref{eq:confidence} against per-particle oracle utility. Every signal is recorded on the \emph{same} captured particles, and the weighting rules in the last row re-weight that same ensemble, so the comparison carries no difference in test-time compute: whatever separates the rules is attributable to the coefficients alone. The capture uses $M{=}8$, since $\hat D_m$ is identically zero for $M\le2$, and the diagnostic forces every coefficient to be non-degenerate even when the run's own configuration set one to zero, so that the rules stay comparable. Note that $\bar c_m$ is recoverable from any existing capture with no re-decoding, since $\bar c_m$ equals the stored average log-probability plus the stored mean entropy; only the varentropy $V_u$ of Eq.~\eqref{eq:varentropy} requires a re-capture pass.

Two questions are answered in turn. \textbf{Is disagreement the better per-particle signal?} The first two rows plot each particle's oracle utility against $D_m$ and against $\bar H_m$. Utility here is $\operatorname{Recall}@B$ against the oracle score $S^\star$ (Eq.~\eqref{eq:recall_at_b}), i.e.\ measured on the retention a single particle induces rather than on its queries alone; the capture also records each particle's standalone coverage of $\mathcal{Q}_{\mathrm{future}}$ (Eq.~\eqref{eq:coverage}) and the two orderings agree, so the figure reports the downstream quantity. A negative slope in the first row is the premise of the penalty: a trajectory the rest of the ensemble does not reproduce keeps tokens the true future queries would not have kept. The third row ranks every candidate signal by its Spearman correlation with the same utility, sign-oriented so that larger always means ``better particle''. Two comparisons matter there. The first is $\bar c_m$ against the raw $\log p_\theta$: if removing the entropy confound raises the correlation, the reparametrization improves the per-particle signal itself, measured on this paper's own diagnostic. The second is whether $-\hat D_m$ ranks above $-\bar H_m$, which would mean disagreement carries information entropy does not. On LongBench the correlations are \tbd\ for $\log p_\theta$, \tbd\ for $\bar c_m$, \tbd\ for $-\bar H_m$ and \tbd\ for $-\hat D_m$, with \tbd, \tbd, \tbd\ and \tbd\ on RULER. Because a marginal correlation cannot separate a signal from what it shares with the others, the \emph{semipartial} correlation of $D_m$ with the residual of utility on $(\bar c_m,\bar H_m)$ is what gates $\lambda_D>0$ in the end-to-end grid.

\textbf{Does the weighting built on it win?} The last row reports ensemble coverage of $\mathcal{Q}_{\mathrm{future}}$ under the $(\beta,\gamma,\lambda_D)$ cells $(0,0,0)$, $(1,1,0)$, $(1,0,0)$, $(1,1.3,0)$, $(1,1,\lambda_D)$ and $(1,0,\lambda_D)$---i.e.\ including the uniform and no-entropy controls that the earlier parametrization could not express. Consistent with the calibration result of Figure~\ref{fig:confidence_calibration}---the reliability signal is real but weak---the absolute differences between the bars are small, which is exactly why the weighting stays soft and $\lambda_D$ is bounded. The statistic to read is therefore not the bar height but the \emph{paired} per-sample comparison printed in the panel, since absolute coverage varies far more across prompts than across weighting rules: the disagreement rule improves mean coverage by \tbd\ and wins on \tbd\% of samples on LongBench, and by \tbd\ on \tbd\% on RULER. This score-level evidence is available for both variants now; the corresponding end-to-end accuracies are the \tbd\ rows of Tables~\ref{tab:longbench_mistral}, \ref{tab:ruler_main}, \ref{tab:longbench_llama}, \ref{tab:ablation_summary}, \ref{tab:ruler_ablation} and~\ref{tab:ablation_summary}.

\begin{figure}[htbp]
  \centering
  \includegraphics[width=\textwidth]{figures_combined/fig10_particle_divergence.pdf}
  \caption{\textbf{Particle reliability signals.} Oracle retention utility versus trajectory disagreement and predictive entropy, signal correlations, and ensemble coverage under alternative weighting rules. All rules use identical particles. Mistral-7B-Instruct-v0.2, $B=128$, $T=8$, $M=8$.}
  \label{fig:particle_divergence}
\end{figure}


\subsection{Retained-Token Composition}
\label{app:composition}

Figure~\ref{fig:retained_token_composition} assigns each prompt token, in priority order, to attention sink, recency window, answer/evidence, number, punctuation/boilerplate, capitalized entity, or other. It reports both the fraction of heads retaining a token and the enrichment
$P(c\mid\mathcal R)/P(c\mid X)$.

The protected sink and recency classes are identical across methods. On LongBench, answer/evidence enrichment is $1.03$ for attention, $1.00$ for the AnDPro window, and $2.51$ for \lorekvE; the corresponding retention rates are approximately $0.028$, $0.028$, and $0.069$. On RULER, answer/evidence enrichment is $16.5$, $16.3$, and $20.3$, respectively. Other token classes change only modestly, indicating that the main compositional difference is increased retention of answer-bearing evidence. Composition is a property of the retained set, so the \lorekvD\ counterpart is obtained by re-weighting the same captured particles under the disagreement rule and re-running selection only, with no additional decoding; it is not yet reported. This complements the retrieval analysis of Section~\ref{subsec:niah}.

\begin{figure}[t]
  \centering
  \includegraphics[width=\textwidth]{figures_combined/fig7_retained_token_composition.pdf}
  \caption{\textbf{Retention by token class.} Mean retained fraction (left) and class enrichment $P(c\mid\mathcal R)/P(c\mid X)$ (right).}
  \label{fig:retained_token_composition}
\end{figure}


\subsection{Layer- and Head-Level Structure}
\label{app:layer_head}

Figure~\ref{fig:layer_head_analysis} localizes the perturbation signal across the network, under the \lorekvE\ weights of Appendix~\ref{app:mechanism_protocol}. Layer perturbation energy increases with depth, while the gain
$G_{l,h}$ from replacing window-only queries with multi-lookahead queries is positive for most heads but strongly non-uniform. The largest gains concentrate in a subset of middle-depth heads. This heterogeneity is consistent with the head specialization motivating head-wise retention methods \citep{feng2024adakv,fu2024headkv} and provides the premise for the adaptive allocation rule in Eq.~\eqref{eq:adaptive_budget}; its small end-to-end ablation gain indicates that most of the improvement is already captured by the core token score.

\begin{figure}[t]
  \centering
  \includegraphics[width=\textwidth]{figures_combined/fig8_layer_head_analysis.pdf}
  \caption{\textbf{Future-query signal across layers and heads.} Layer perturbation energy (left) and retention-recall gain from multi-lookahead queries (right). Mistral-7B-Instruct-v0.2, $B=128$, $M=4$.}
  \label{fig:layer_head_analysis}
\end{figure}


\subsection{Budget Equivalence in Behavioral Fidelity}
\label{app:budget_curve}

Figure~\ref{fig:budget_equivalence} complements Table~\ref{tab:budget_equivalence} by measuring cache quality continuously with the behavioral fidelity $F(B)$ of Eq.~\eqref{eq:fidelity}. On LongBench, \lorekvE\ increases from $0.852$ at $B{=}64$ to $0.981$ at $B{=}1024$, compared with $0.818\!\rightarrow\!0.975$ for attention and $0.818\!\rightarrow\!0.973$ for the AnDPro window. RULER shows the same ordering over a narrower range: $0.916\!\rightarrow\!0.989$ for \lorekvE\ versus $0.900\!\rightarrow\!0.986$ and $0.902\!\rightarrow\!0.987$ for the two prompt-side criteria.

Matching \lorekvE\ fidelity requires approximately $33$--$43\%$ more retained-budget reduction on LongBench and $27$--$34\%$ on RULER across the measured range. These values concern behavioral fidelity rather than task accuracy and therefore complement, rather than replace, Table~\ref{tab:budget_equivalence}. Where a baseline never reaches the target fidelity within the sweep, the corresponding saving is only lower-bounded.

\begin{figure}[t]
  \centering
  \includegraphics[width=\textwidth]{figures_combined/fig9_budget_equivalence.pdf}
  \caption{\textbf{Budget equivalence under behavioral fidelity.} Fidelity across cache budgets (left) and baseline cache savings required to match \lorekvE\ (right).}
  \label{fig:budget_equivalence}
\end{figure}

\section{The Two Reliability Variants in the $(\beta,\gamma,\lambda_D)$ Plane} \label{app:reparam_grid}

The paper carries two variants end to end, as points of the reparametrized plane: \lorekvE\ at $(\beta,\gamma,\lambda_D)=(1,1.3,0)$, whose LongBench averages with Mistral-7B-Instruct-v0.2 are $40.28$, $41.38$, $41.84$ and $42.21$ at $B=128$, $256$, $512$ and $1024$; and \lorekvD\ at $(1,1,0.3)$, whose end-to-end runs are \tbd. The reparametrization is what makes the pair well-posed at all---in the earlier $(\lambda_H,\eta)$ parametrization the two settings were not comparable coefficient-for-coefficient, and $\lambda_H{=}0$ silently carried a full unit of entropy penalty (Proposition~1). Read together, \lorekvD\ is not \lorekvE\ with the entropy channel swapped out: it keeps that channel at $\gamma{=}1$ and adds disagreement beside it, so the pair differs in \emph{two} coefficients and is a variant comparison rather than a single-factor ablation. The single-factor cells that isolate one coefficient at a time---uniform $(0,0,0)$, pure typicality $(1,0,0)$, disagreement-without-entropy $(1,0,0.3)$ and stronger entropy $(1,2,0)$---are ablation rows of Table~\ref{tab:ablation_summary}, not variants of the method, and are reported there. \lorekvD\ is calibrated per Section~\ref{subsec:confidence} so that $\lambda_D\hat D_m$ and $\gamma\bar H_m$ induce comparable realized spread.

\textbf{Protocol.} The comparison runs at $B{=}128$ only: Table~\ref{tab:ruler_main} shows the criterion separating by $5.05$ points at $B{=}128$ and $0.42$ at $B{=}1024$, so the larger budgets have no resolving power for a weighting rule. Both variants re-weight nested prefixes of a \emph{single} eight-particle draw per prompt, so neither differs from the other by test-time compute. We run three seeds minimum: the gaps at stake are $0.01$ in Table~\ref{tab:ablation_summary} and $0.04$ in Table~\ref{tab:ruler_ablation}, and the existing $M$-sweep is non-monotone ($M{=}2$ at $41.10$/$41.22$ sits below $M{=}1$ at $41.56$/$42.03$), so the noise floor is unknown and may well exceed them. If it does, ``\lorekvD\ does not separate from \lorekvE'' is the correct result and is what we will report.


\section{Limitations} \label{app:limitations}
LORE-KV trades one-time pre-generation computation for a better future-query reference. The lookahead trajectories are generated before decoding, preserving the same cache size and per-token decoding cost as other eviction methods, but adding noticeable wall-clock overhead. Unlike learned methods such as LookaheadKV, LORE-KV requires no offline training or model-specific learned parameters.

\textcolor{red}{LORE-KV decides which KV pairs should be kept, but once a KV pair is removed, its information is lost. In contrast, methods such as IndexMem keep a compressed summary of evicted information \citep{yang2026indexmem}, while LaProx considers interactions across heads and layers when selecting what to retain \citep{mai2026laprox}. LORE-KV also compresses the cache only once after the prompt, so extending it to update the cache during long or agentic generations is left for future work.}

More broadly, future work should study how LORE-KV can preserve important instructions, structural information, and task-relevant evidence under aggressive compression, especially in multi-instruction, safety-sensitive, and long-context settings.
\end{document}